\documentclass[output=paper,colorlinks,citecolor=brown
% ,hidelinks
% showindex
]{langscibook}
\ChapterDOI{10.5281/zenodo.20285188}

\author{Nadja Fiebig{}\affiliation{Leipzig University}}
\title{Two types of verbal number in Kalenjin pluractionality}
\abstract{This paper investigates pluractionality in Pokot and Kipsigis, two closely related varieties of Kalenjin, a macrolanguage belonging to the Southern Nilotic branch of Nilo-Saharan \citep{Ethnologue2023}.  Most typological literature (e.g. \citealt{Corbett2000}), as well as generative analyses (\citealt{Amato2018}, \citealt{Thornton2019}), treat participant number and pluractionality (=event number) as one category, called verbal number. However, I show that these two categories cannot be conflated in Kalenjin by providing novel data, coming from original fieldwork conducted in Kenya. I further provide the first morphological and semantic description of pluractionality in both varieties.}


\IfFileExists{../localcommands.tex}{
   \addbibresource{../localbibliography.bib}
   %%%%%%%%%%%%%%%%%%%%%%%%%%%%%%%%%%%%%%%%%%%%%
%%%%%%%%%% From LSP %%%%%%%%%%%%%%%%%%%%%%%%%
%%%%%%%%%%%%%%%%%%%%%%%%%%%%%%%%%%%%%%%%%%%%%

\usepackage{tabularx,multicol}
\usepackage{url}
\urlstyle{same}

\usepackage{langsci-optional}
\usepackage{langsci-lgr}
\usepackage{langsci-gb4e}

% NOTE THAT THE FOLLOWING PACKAGES ARE BANNED: 
% - pstricks
% - tipa
%%%%%%%%%%%%%%%%%%%%%%%%%%%%%%%%%%%%%%%%%%%%%
%%%%%%%%%% From ACAL LaTeX template %%%%%%%%%
%%%%%%%%%%%%%%%%%%%%%%%%%%%%%%%%%%%%%%%%%%%%%

%For stacking text, used here in autosegmental diagrams
\usepackage{stackengine}

%To combine rows in tables
\usepackage{multirow}

%Gives some extra formatting options, e.g. underlining/strikeout
\usepackage[normalem]{ulem}

%Use package/command below to create a double-spaced document, if you want one. Uncomment BOTH the package and the command (\doublespacing) to create a doublespaced document, or leave them as is to have a single-spaced document.
%\usepackage{setspace}
%\doublespacing 

 

%use for special OT tableaux symbols like bomb and sad face. must be loaded early on because it doesn't play well with some other packages
\usepackage{fourier-orns}

%Basic math symbols 
\usepackage{pifont}
\usepackage{amssymb}

%%%Gives shortcuts for glossing. The use of this package is NOT explained in the Quick Reference Guide, but the documentation is on CTAN for those that are interested. MJKD finds it handy for glossing. (https://ctan.org/pkg/leipzig?lang=en)
\usepackage{leipzig}

%Tables
% \usepackage{caption} %For table captions 

%For OT-style tableaux
\usepackage[notipa]{ot-tableau}

%highlights text with \hl{text}
\usepackage{color, soul} %note that soul sometimes eats Unicode text witouth telling you. 

%Drawing Syntax Trees
\usepackage[linguistics]{forest}

%This specifies some formatting for the forest trees to make them nicer to look at. Unfortunately this was borrowed by Michael Diercks from someone a while ago and he can't remember who. Apologies and if someone lets him know he will update this.
\forestset{
  nice nodes/.style={
    for tree={
      inner sep=0pt,
      fit=band,
    },
  },
  default preamble=nice nodes,
}

%A couple of packages that seemed to prefer being called toward the end of the preamble

%This package provides macros for typesetting SPE-style phonological rules. I've redefined the \oneof command here, so that there the rule includes a right brace. 
\usepackage{phonrule}
\renewcommand{\oneof}[2][c]{\ensuremath{
    ~\left\{
    \begin{tabular}{#1#1}#2\end{tabular}
    \right\}
    }}

%For using Greek letters outside of math mode.
% \usepackage{textgreek}


%Random, lets us use the XeLaTeX logo. Not important to the template at all.
% \usepackage{metalogo}



%%%%%%%%%%%%%%%%%%%%%%%%%%%%%%%%%%%%%%%%%%%%%
%%%%%%%%%% From Smith paper %%%%%%%%%%%%%%%%%
%%%%%%%%%%%%%%%%%%%%%%%%%%%%%%%%%%%%%%%%%%%%%

\usetikzlibrary{positioning}
\usetikzlibrary{trees}
% \usepackage{pstricks} %pstricks is banned. Use tikz instead. 
% \usepackage{pst-node}
%\usepackage{tikz}

%%%%%%%%%%%%%%%%%%%%%%%%%%%%%%%%%%%%%%%%%%%%%
%%%%%%%%%% From Gluckman paper %%%%%%%%%%%%%%
%%%%%%%%%%%%%%%%%%%%%%%%%%%%%%%%%%%%%%%%%%%%%

\usepackage{amsmath}



%%%%%%%%%%%%%%%%%%%%%%%%%%%%%%%%%%%%%%%%%%%%%
%%%%%%%%%% From Kandybowicz paper %%%%%%%%%%%
%%%%%%%%%%%%%%%%%%%%%%%%%%%%%%%%%%%%%%%%%%%%%

%These are in the .tex, but there's nothing in the "Figures" folder so I'm not sure that it's actually doing anything. 
 
% \graphicspath{ {./Figures/} }

%%%%%%%%%%%%%%%%%%%%%%%%%%%%%%%%%%%%%%%%%%%%%
%%%%%%%%%% From Matte paper %%%%%%%%%%%%%%%%%
%%%%%%%%%%%%%%%%%%%%%%%%%%%%%%%%%%%%%%%%%%%%%

%%%%%%%%%%%%%%%%%%%%%%%%%%%%%%%%%%%%%%%%%%%%%
%%%%%%%%%% From Grollemund paper %%%%%%%%%%%%%%
%%%%%%%%%%%%%%%%%%%%%%%%%%%%%%%%%%%%%%%%%%%%%

% \usepackage{lscape}

%%%%%%%%%%%%%%%%%%%%%%%%%%%%%%%%%%%%%%%%%%%%%
%%%%%%%%%% From Burns paper %%%%%%%%%%%%%%
%%%%%%%%%%%%%%%%%%%%%%%%%%%%%%%%%%%%%%%%%%%%%


% \usepackage{url}
% \urlstyle{same}

\usepackage{listings}
\lstset{basicstyle=\ttfamily,tabsize=2,breaklines=true}

%MJKD commented out - these are standard for chapters in edited volumes, I don't think we need them here in the preamble. 

% \usepackage{tabularx,multicol}
% \usepackage{langsci-basic}
% \usepackage{langsci-gb4e}
% \bibliography{localbibliography}
% \newcommand{\orcid}[1]{}
% \pagenumbering{roman}

% \IfFileExists{../localcommands.tex}{%hack to check whether this is being compiled as part of a collection or standalone
%    %%%%%%%%%%%%%%%%%%%%%%%%%%%%%%%%%%%%%%%%%%%%%
%%%%%%%%%% From LSP %%%%%%%%%%%%%%%%%%%%%%%%%
%%%%%%%%%%%%%%%%%%%%%%%%%%%%%%%%%%%%%%%%%%%%%

\usepackage{tabularx,multicol}
\usepackage{url}
\urlstyle{same}

\usepackage{langsci-optional}
\usepackage{langsci-lgr}
\usepackage{langsci-gb4e}

% NOTE THAT THE FOLLOWING PACKAGES ARE BANNED: 
% - pstricks
% - tipa
%%%%%%%%%%%%%%%%%%%%%%%%%%%%%%%%%%%%%%%%%%%%%
%%%%%%%%%% From ACAL LaTeX template %%%%%%%%%
%%%%%%%%%%%%%%%%%%%%%%%%%%%%%%%%%%%%%%%%%%%%%

%For stacking text, used here in autosegmental diagrams
\usepackage{stackengine}

%To combine rows in tables
\usepackage{multirow}

%Gives some extra formatting options, e.g. underlining/strikeout
\usepackage[normalem]{ulem}

%Use package/command below to create a double-spaced document, if you want one. Uncomment BOTH the package and the command (\doublespacing) to create a doublespaced document, or leave them as is to have a single-spaced document.
%\usepackage{setspace}
%\doublespacing 

 

%use for special OT tableaux symbols like bomb and sad face. must be loaded early on because it doesn't play well with some other packages
\usepackage{fourier-orns}

%Basic math symbols 
\usepackage{pifont}
\usepackage{amssymb}

%%%Gives shortcuts for glossing. The use of this package is NOT explained in the Quick Reference Guide, but the documentation is on CTAN for those that are interested. MJKD finds it handy for glossing. (https://ctan.org/pkg/leipzig?lang=en)
\usepackage{leipzig}

%Tables
% \usepackage{caption} %For table captions 

%For OT-style tableaux
\usepackage[notipa]{ot-tableau}

%highlights text with \hl{text}
\usepackage{color, soul} %note that soul sometimes eats Unicode text witouth telling you. 

%Drawing Syntax Trees
\usepackage[linguistics]{forest}

%This specifies some formatting for the forest trees to make them nicer to look at. Unfortunately this was borrowed by Michael Diercks from someone a while ago and he can't remember who. Apologies and if someone lets him know he will update this.
\forestset{
  nice nodes/.style={
    for tree={
      inner sep=0pt,
      fit=band,
    },
  },
  default preamble=nice nodes,
}

%A couple of packages that seemed to prefer being called toward the end of the preamble

%This package provides macros for typesetting SPE-style phonological rules. I've redefined the \oneof command here, so that there the rule includes a right brace. 
\usepackage{phonrule}
\renewcommand{\oneof}[2][c]{\ensuremath{
    ~\left\{
    \begin{tabular}{#1#1}#2\end{tabular}
    \right\}
    }}

%For using Greek letters outside of math mode.
% \usepackage{textgreek}


%Random, lets us use the XeLaTeX logo. Not important to the template at all.
% \usepackage{metalogo}



%%%%%%%%%%%%%%%%%%%%%%%%%%%%%%%%%%%%%%%%%%%%%
%%%%%%%%%% From Smith paper %%%%%%%%%%%%%%%%%
%%%%%%%%%%%%%%%%%%%%%%%%%%%%%%%%%%%%%%%%%%%%%

\usetikzlibrary{positioning}
\usetikzlibrary{trees}
% \usepackage{pstricks} %pstricks is banned. Use tikz instead. 
% \usepackage{pst-node}
%\usepackage{tikz}

%%%%%%%%%%%%%%%%%%%%%%%%%%%%%%%%%%%%%%%%%%%%%
%%%%%%%%%% From Gluckman paper %%%%%%%%%%%%%%
%%%%%%%%%%%%%%%%%%%%%%%%%%%%%%%%%%%%%%%%%%%%%

\usepackage{amsmath}



%%%%%%%%%%%%%%%%%%%%%%%%%%%%%%%%%%%%%%%%%%%%%
%%%%%%%%%% From Kandybowicz paper %%%%%%%%%%%
%%%%%%%%%%%%%%%%%%%%%%%%%%%%%%%%%%%%%%%%%%%%%

%These are in the .tex, but there's nothing in the "Figures" folder so I'm not sure that it's actually doing anything. 
 
% \graphicspath{ {./Figures/} }

%%%%%%%%%%%%%%%%%%%%%%%%%%%%%%%%%%%%%%%%%%%%%
%%%%%%%%%% From Matte paper %%%%%%%%%%%%%%%%%
%%%%%%%%%%%%%%%%%%%%%%%%%%%%%%%%%%%%%%%%%%%%%

%%%%%%%%%%%%%%%%%%%%%%%%%%%%%%%%%%%%%%%%%%%%%
%%%%%%%%%% From Grollemund paper %%%%%%%%%%%%%%
%%%%%%%%%%%%%%%%%%%%%%%%%%%%%%%%%%%%%%%%%%%%%

% \usepackage{lscape}

%%%%%%%%%%%%%%%%%%%%%%%%%%%%%%%%%%%%%%%%%%%%%
%%%%%%%%%% From Burns paper %%%%%%%%%%%%%%
%%%%%%%%%%%%%%%%%%%%%%%%%%%%%%%%%%%%%%%%%%%%%


% \usepackage{url}
% \urlstyle{same}

\usepackage{listings}
\lstset{basicstyle=\ttfamily,tabsize=2,breaklines=true}

%MJKD commented out - these are standard for chapters in edited volumes, I don't think we need them here in the preamble. 

% \usepackage{tabularx,multicol}
% \usepackage{langsci-basic}
% \usepackage{langsci-gb4e}
% \bibliography{localbibliography}
% \newcommand{\orcid}[1]{}
% \pagenumbering{roman}

% \IfFileExists{../localcommands.tex}{%hack to check whether this is being compiled as part of a collection or standalone
%    %%%%%%%%%%%%%%%%%%%%%%%%%%%%%%%%%%%%%%%%%%%%%
%%%%%%%%%% From LSP %%%%%%%%%%%%%%%%%%%%%%%%%
%%%%%%%%%%%%%%%%%%%%%%%%%%%%%%%%%%%%%%%%%%%%%

\usepackage{tabularx,multicol}
\usepackage{url}
\urlstyle{same}

\usepackage{langsci-optional}
\usepackage{langsci-lgr}
\usepackage{langsci-gb4e}

% NOTE THAT THE FOLLOWING PACKAGES ARE BANNED: 
% - pstricks
% - tipa
%%%%%%%%%%%%%%%%%%%%%%%%%%%%%%%%%%%%%%%%%%%%%
%%%%%%%%%% From ACAL LaTeX template %%%%%%%%%
%%%%%%%%%%%%%%%%%%%%%%%%%%%%%%%%%%%%%%%%%%%%%

%For stacking text, used here in autosegmental diagrams
\usepackage{stackengine}

%To combine rows in tables
\usepackage{multirow}

%Gives some extra formatting options, e.g. underlining/strikeout
\usepackage[normalem]{ulem}

%Use package/command below to create a double-spaced document, if you want one. Uncomment BOTH the package and the command (\doublespacing) to create a doublespaced document, or leave them as is to have a single-spaced document.
%\usepackage{setspace}
%\doublespacing 

 

%use for special OT tableaux symbols like bomb and sad face. must be loaded early on because it doesn't play well with some other packages
\usepackage{fourier-orns}

%Basic math symbols 
\usepackage{pifont}
\usepackage{amssymb}

%%%Gives shortcuts for glossing. The use of this package is NOT explained in the Quick Reference Guide, but the documentation is on CTAN for those that are interested. MJKD finds it handy for glossing. (https://ctan.org/pkg/leipzig?lang=en)
\usepackage{leipzig}

%Tables
% \usepackage{caption} %For table captions 

%For OT-style tableaux
\usepackage[notipa]{ot-tableau}

%highlights text with \hl{text}
\usepackage{color, soul} %note that soul sometimes eats Unicode text witouth telling you. 

%Drawing Syntax Trees
\usepackage[linguistics]{forest}

%This specifies some formatting for the forest trees to make them nicer to look at. Unfortunately this was borrowed by Michael Diercks from someone a while ago and he can't remember who. Apologies and if someone lets him know he will update this.
\forestset{
  nice nodes/.style={
    for tree={
      inner sep=0pt,
      fit=band,
    },
  },
  default preamble=nice nodes,
}

%A couple of packages that seemed to prefer being called toward the end of the preamble

%This package provides macros for typesetting SPE-style phonological rules. I've redefined the \oneof command here, so that there the rule includes a right brace. 
\usepackage{phonrule}
\renewcommand{\oneof}[2][c]{\ensuremath{
    ~\left\{
    \begin{tabular}{#1#1}#2\end{tabular}
    \right\}
    }}

%For using Greek letters outside of math mode.
% \usepackage{textgreek}


%Random, lets us use the XeLaTeX logo. Not important to the template at all.
% \usepackage{metalogo}



%%%%%%%%%%%%%%%%%%%%%%%%%%%%%%%%%%%%%%%%%%%%%
%%%%%%%%%% From Smith paper %%%%%%%%%%%%%%%%%
%%%%%%%%%%%%%%%%%%%%%%%%%%%%%%%%%%%%%%%%%%%%%

\usetikzlibrary{positioning}
\usetikzlibrary{trees}
% \usepackage{pstricks} %pstricks is banned. Use tikz instead. 
% \usepackage{pst-node}
%\usepackage{tikz}

%%%%%%%%%%%%%%%%%%%%%%%%%%%%%%%%%%%%%%%%%%%%%
%%%%%%%%%% From Gluckman paper %%%%%%%%%%%%%%
%%%%%%%%%%%%%%%%%%%%%%%%%%%%%%%%%%%%%%%%%%%%%

\usepackage{amsmath}



%%%%%%%%%%%%%%%%%%%%%%%%%%%%%%%%%%%%%%%%%%%%%
%%%%%%%%%% From Kandybowicz paper %%%%%%%%%%%
%%%%%%%%%%%%%%%%%%%%%%%%%%%%%%%%%%%%%%%%%%%%%

%These are in the .tex, but there's nothing in the "Figures" folder so I'm not sure that it's actually doing anything. 
 
% \graphicspath{ {./Figures/} }

%%%%%%%%%%%%%%%%%%%%%%%%%%%%%%%%%%%%%%%%%%%%%
%%%%%%%%%% From Matte paper %%%%%%%%%%%%%%%%%
%%%%%%%%%%%%%%%%%%%%%%%%%%%%%%%%%%%%%%%%%%%%%

%%%%%%%%%%%%%%%%%%%%%%%%%%%%%%%%%%%%%%%%%%%%%
%%%%%%%%%% From Grollemund paper %%%%%%%%%%%%%%
%%%%%%%%%%%%%%%%%%%%%%%%%%%%%%%%%%%%%%%%%%%%%

% \usepackage{lscape}

%%%%%%%%%%%%%%%%%%%%%%%%%%%%%%%%%%%%%%%%%%%%%
%%%%%%%%%% From Burns paper %%%%%%%%%%%%%%
%%%%%%%%%%%%%%%%%%%%%%%%%%%%%%%%%%%%%%%%%%%%%


% \usepackage{url}
% \urlstyle{same}

\usepackage{listings}
\lstset{basicstyle=\ttfamily,tabsize=2,breaklines=true}

%MJKD commented out - these are standard for chapters in edited volumes, I don't think we need them here in the preamble. 

% \usepackage{tabularx,multicol}
% \usepackage{langsci-basic}
% \usepackage{langsci-gb4e}
% \bibliography{localbibliography}
% \newcommand{\orcid}[1]{}
% \pagenumbering{roman}

% \IfFileExists{../localcommands.tex}{%hack to check whether this is being compiled as part of a collection or standalone
%    \input{../localpackages}
%    \input{../localcommands}
%    \input{../localhyphenation}
%     \bibliography{localbibliography}
%     \togglepaper[23]
% }{}

% %custom footer for preprints
% \papernote{\scriptsize\normalfont
%     Change author.
%     Change title.
%     To appear in:
%     Change Volume Editor.  
%     Change volume title.  
%     Berlin: Language Science Press. [preliminary page numbering]
% }



%%%%%%%%%%%%%%%%%%%%%%%%%%%%%%%%%%%%%%%%%%%%%
%%%%%%%%%% From Feibig paper %%%%%%%%%%%%%%
%%%%%%%%%%%%%%%%%%%%%%%%%%%%%%%%%%%%%%%%%%%%%

\usepackage{tikz-qtree}
\usepackage{tikz-qtree-compat}

%%%%%%%%%%%%%%%%%%%%%%%%%%%%%%%%%%%%%%%%%%%%%
%%%%%%%%%% From Olson paper %%%%%%%%%%%%%%
%%%%%%%%%%%%%%%%%%%%%%%%%%%%%%%%%%%%%%%%%%%%%

%MJKD - TIPA is forbidden lol. Commenting out for now but we'll have to address this in the Olson paper

%\usepackage[tone]{tipa}

%%%%%%%%%%%%%%%%%%%%%%%%%%%%%%%%%%%%%%%%%%%%%
%%%%%%%%%% From Caragine paper %%%%%%%%%%%%%%
%%%%%%%%%%%%%%%%%%%%%%%%%%%%%%%%%%%%%%%%%%%%%

\usepackage{placeins}
% \usepackage{graphicx}
% \usepackage{float} %Overleaf suggested this as a solution


%%%%%%%%%%%%%%%%%%%%%%%%%%%%%%%%%%%%%%%%%%%%%
%%%%%%%%%% From Kieffer paper %%%%%%%%%%%%%%
%%%%%%%%%%%%%%%%%%%%%%%%%%%%%%%%%%%%%%%%%%%%%

\usepackage{array}

%%%%%%%%%%%%%%%%%%%%%%%%%%%%%%%%%%%%%%%%%%%%%
%%%%%%%%%% From Payne paper %%%%%%%%%%%%%%
%%%%%%%%%%%%%%%%%%%%%%%%%%%%%%%%%%%%%%%%%%%%%

\usepackage{enumerate}
\usepackage{array,ragged2e}
\usepackage{pgfplots}


%%%%%%%%%%%%%%%%%%%%%%%%%%%%%%%%%%%%%%%%%%%%%
%%%%%%%%%% From Diercks paper %%%%%%%%%%%%%%
%%%%%%%%%%%%%%%%%%%%%%%%%%%%%%%%%%%%%%%%%%%%%

% \usepackage{cgloss}



%%%%%%%%%%%%%%%%%%%%%%%%%%%%%%%%%%%%%%%%%%%%%
%%%%%%%%%% From Li paper %%%%%%%%%%%%%%
%%%%%%%%%%%%%%%%%%%%%%%%%%%%%%%%%%%%%%%%%%%%%



%%%%%%%%%%%%%%%%%%%%%%%%%%%%%%%%%%%%%%%%%%%%%
%%%%%%%%%% From Myers paper %%%%%%%%%%%%%%
%%%%%%%%%%%%%%%%%%%%%%%%%%%%%%%%%%%%%%%%%%%%%



\usepackage{tikz-qtree}
\usepackage{tikz-qtree-compat}


%Package that makes glossing slightly easier.
\RequirePackage{leipzig}
%\RequirePackage{mathtools} % for \mathrlap

% % % Shortcuts (various borrowed from Jason Zentz)
\RequirePackage{xspace}
\xspaceaddexceptions{]\}}
\usepackage{langsci-textipa}
\usepackage{langsci-branding}
\usepackage{tabto}

%    %%%%%%%%%%%%%%%%%%%%%%%%%%%%%%%%%%%%%%%%%%%%%
%%%%%%%%%% From LSP %%%%%%%%%%%%%%%%%%%%%%%%%
%%%%%%%%%%%%%%%%%%%%%%%%%%%%%%%%%%%%%%%%%%%%%


% \newcommand*{\orcid}{}


\makeatletter
\let\thetitle\@title
\let\theauthor\@author
\makeatother

\newcommand{\togglepaper}[1][0]{
   \bibliography{../localbibliography}
   \papernote{\scriptsize\normalfont
     \theauthor.
     \titleTemp.
     To appear in:
     E. Di Tor \& Herr Rausgeberin (ed.).
     Booktitle in localcommands.tex.
     Berlin: Language Science Press. [preliminary page numbering]
   }
   \pagenumbering{roman}
   \setcounter{chapter}{#1}
   \addtocounter{chapter}{-1}
}

\newbool{bookcompile}
\booltrue{bookcompile}
\newcommand{\bookorchapter}[2]{\ifbool{bookcompile}{#1}{#2}}


%%%%%%%%%%%%%%%%%%%%%%%%%%%%%%%%%%%%%%%%%%%%%
%%%%%%%%%% From ACAL LaTeX template %%%%%%%%%
%%%%%%%%%%%%%%%%%%%%%%%%%%%%%%%%%%%%%%%%%%%%%


\usepackage{tikz}

\newcommand{\circled}[1]{\begin{tikzpicture}[baseline=(word.base)]
\node[draw, rounded corners, text height=8pt, text depth=2pt, inner sep=2pt, outer sep=0pt, use as bounding box] (word) {#1};
\end{tikzpicture}
}


%%%%%%%%%%%%%%%%%%%%%%%%%%%%%%%%%%%%%%%%%%%%%%
%%%%%%%%%%%%%%%%%%%%%%%%%%%%%%%%%%%%%%%%%%%%%%

% Useful Ling Shortcuts


%This makes the \emptyset command be a nicer one
\let\oldemptyset\emptyset
\let\emptyset\varnothing
\newcommand{\nothing}{$\emptyset$}

%Not all of these are explained in the Quick Reference Guide, but they are here bc they are relevant to some of our students. 
\newcommand{\xbar}{\ensuremath{'}\xspace}
\newcommand{\xhead}{\ensuremath{^\circ}\xspace}
\newcommand{\Lb}[1]{$\text{[}_{\text{#1}}$ } %A more convenient left bracket
\newcommand{\Rb}[1]{$\text{]}_{\text{#1}}$ } %A more convenient left bracket
\newcommand{\gap}{\underline{\hspace{1.2em}}}
\newcommand{\vP}{\emph{v}P}
\newcommand{\lilv}{\emph{v}}
\newcommand{\Abar}{A$'$-} %A more convenient A-bar notation
\newcommand{\ph}{$\varphi$\xspace} %A more convenient phi
\newcommand{\pro}{\emph{pro}\xspace}
\newcommand{\subs}[1]{\textsubscript{#1}} %A more convenient subscript
\newcommand{\spells}{$\Longleftrightarrow$} %spellout arrow for morph spellout rules
\newcommand{\tr}[1]{\textit{t}\textsubscript{\textit{#1}}} %easy traces with subscript
\newcommand{\supers}[1]{\textsuperscript{#1}}

%These are borrowed/modified from Andrew Mackenzie's ling-macros package. We have copied them in here (instead of just using the package itself) to avoid a minor clash that was generating an error.


% Abbreviations for glossing, based on Leipzig
\newleipzig{hab}{hab}{habitual}
\newleipzig{rem}{rem}{remote}
\newleipzig{sm}{sm}{subject marker}
\newleipzig{t}{t}{tense}
\newleipzig{aa}{aa}{anti-agreement}
\newleipzig{pron}{pron}{pronoun}
\newleipzig{rec}{rec}{recent}
\newleipzig{om}{om}{object marker}
%\newleipzig{ipfv}{ipfv}{imperfective}
\newleipzig{asp}{asp}{aspect}
\newleipzig{lk}{lk}{linker}
\newleipzig{pcl}{pcl}{particle}
\newleipzig{stat}{stat}{stative}
\newleipzig{ints}{ints}{intensive}
\newleipzig{ascl}{ascl}{assertive subject clitic}
\newleipzig{nascl}{nascl}{non-assertive subject clitic}
\newleipzig{ta}{ta}{tense and/or aspect}
\newleipzig{assoc}{assoc}{associative marker}
\newleipzig{hon}{hon}{honorific}
%\newleipzig{whprt}{wh}{\wh particle}
\newleipzig{sa}{sa}{subject agreement}
\newleipzig{conj}{conj}{conjunction}
%\newleipzig{loc}{loc}{locative}
\newleipzig{expl}{expl}{expletive}
\newleipzig{rcm}{rcm}{reciprocal marker}
\newleipzig{pers}{pers}{persistive}
\newleipzig{fv}{fv}{final vowel}
%\newleipzig{}{}{} %this is just to copy for when I want to add more


%%%%%%%%%%%%%%%%%%%%%%%%%%%%%%%%%%%%%%%%%%%%%%%%%%%%%
%%%%%%%%%%%%%%%%%%%%%%%%%%%%%%%%%%%%%%%%%%%%%%%%%%%%%


%%%%%%%%%%%%%%%%%%%%%%%%%%%%%%%%%%%%%%%%%%%%%
%%%%%%%%%% From Gluckman paper %%%%%%%%%%%%%%
%%%%%%%%%%%%%%%%%%%%%%%%%%%%%%%%%%%%%%%%%%%%%

\newcommand{\pcite}[1]{\citeauthor{#1}'s (\citeyear{#1})}


%%%%%%%%%%%%%%%%%%%%%%%%%%%%%%%%%%%%%%%%%%%%%
%%%%%%%%%% From Myers paper %%%%%%%%%%%%%%
%%%%%%%%%%%%%%%%%%%%%%%%%%%%%%%%%%%%%%%%%%%%%

%% hspace over width of something without showing it
\newcommand*{\hspaceThis}[1]{\settowidth{\LSPTmp}{#1}\hspace*{\LSPTmp}}
% use \hphantom{} for this


%%%%%%%%%%%%%%%%%%%%%%%%%%%%%%%%%%%%%%%%%%%%%
%%%%%%%%%% From Matte paper %%%%%%%%%%%%%%%%%
%%%%%%%%%%%%%%%%%%%%%%%%%%%%%%%%%%%%%%%%%%%%%

%mjkd commented this out in case it's breaking something right now. 
% \newcounter{xlistex}
% \newcommand{\footnoteex}{\stepcounter{xlistex}(\roman{xlistex})}

\newleipzig{sp}{sp}{speaker} %a new leipzig glossing command


%%%%%%%%%%%%%%%%%%%%%%%%%%%%%%%%%%%%%%%%%%%%%
%%%%%%%%%% From GamFranich paper %%%%%%%%%%%%%%
%%%%%%%%%%%%%%%%%%%%%%%%%%%%%%%%%%%%%%%%%%%%%

%MJKD commented these out - both will be provided by the overall LSP book template 

% \bibliography{localbibliography}
% \newcommand{\orcid}[1]{}

%%%%%%%%%%%%%%%%%%%%%%%%%%%%%%%%%%%%%%%%%%%%%
%%%%%%%%%% From Diercks paper %%%%%%%%%%%%%%
%%%%%%%%%%%%%%%%%%%%%%%%%%%%%%%%%%%%%%%%%%%%%

\newcommand{\abar}{$\bar{\text{A}}$\xspace} % this one is different to the \Abar command in Mike's shortcuts doc
\newcommand{\topFeat}{[\textsc{top}]\xspace}


%%%%%%%%%%%%%%%%%%%%%%%%%%%%%%%%%%%%%%%%%%%%%
%%%%%%%%%% From Kinchsular paper %%%%%%%%%%%%%%
%%%%%%%%%%%%%%%%%%%%%%%%%%%%%%%%%%%%%%%%%%%%%

%%%%%%%%%%%%%%%%%%%%%%%%%%%%%%%%%%%%%%%%%%%%%
%%%%%%%%%% From Chen paper %%%%%%%%%%%%%%
%%%%%%%%%%%%%%%%%%%%%%%%%%%%%%%%%%%%%%%%%%%%%

\newcounter{savedex}
\makeatletter
\newcommand*{\saveex}{\setcounter{savedex}{\value{\@xnumctr}}}
\newcommand*{\resumeex}{\setcounter{\@xnumctr}{\value{savedex}}}
\makeatother


%    \input{../localhyphenation}
%     \bibliography{localbibliography}
%     \togglepaper[23]
% }{}

% %custom footer for preprints
% \papernote{\scriptsize\normalfont
%     Change author.
%     Change title.
%     To appear in:
%     Change Volume Editor.  
%     Change volume title.  
%     Berlin: Language Science Press. [preliminary page numbering]
% }



%%%%%%%%%%%%%%%%%%%%%%%%%%%%%%%%%%%%%%%%%%%%%
%%%%%%%%%% From Feibig paper %%%%%%%%%%%%%%
%%%%%%%%%%%%%%%%%%%%%%%%%%%%%%%%%%%%%%%%%%%%%

\usepackage{tikz-qtree}
\usepackage{tikz-qtree-compat}

%%%%%%%%%%%%%%%%%%%%%%%%%%%%%%%%%%%%%%%%%%%%%
%%%%%%%%%% From Olson paper %%%%%%%%%%%%%%
%%%%%%%%%%%%%%%%%%%%%%%%%%%%%%%%%%%%%%%%%%%%%

%MJKD - TIPA is forbidden lol. Commenting out for now but we'll have to address this in the Olson paper

%\usepackage[tone]{tipa}

%%%%%%%%%%%%%%%%%%%%%%%%%%%%%%%%%%%%%%%%%%%%%
%%%%%%%%%% From Caragine paper %%%%%%%%%%%%%%
%%%%%%%%%%%%%%%%%%%%%%%%%%%%%%%%%%%%%%%%%%%%%

\usepackage{placeins}
% \usepackage{graphicx}
% \usepackage{float} %Overleaf suggested this as a solution


%%%%%%%%%%%%%%%%%%%%%%%%%%%%%%%%%%%%%%%%%%%%%
%%%%%%%%%% From Kieffer paper %%%%%%%%%%%%%%
%%%%%%%%%%%%%%%%%%%%%%%%%%%%%%%%%%%%%%%%%%%%%

\usepackage{array}

%%%%%%%%%%%%%%%%%%%%%%%%%%%%%%%%%%%%%%%%%%%%%
%%%%%%%%%% From Payne paper %%%%%%%%%%%%%%
%%%%%%%%%%%%%%%%%%%%%%%%%%%%%%%%%%%%%%%%%%%%%

\usepackage{enumerate}
\usepackage{array,ragged2e}
\usepackage{pgfplots}


%%%%%%%%%%%%%%%%%%%%%%%%%%%%%%%%%%%%%%%%%%%%%
%%%%%%%%%% From Diercks paper %%%%%%%%%%%%%%
%%%%%%%%%%%%%%%%%%%%%%%%%%%%%%%%%%%%%%%%%%%%%

% \usepackage{cgloss}



%%%%%%%%%%%%%%%%%%%%%%%%%%%%%%%%%%%%%%%%%%%%%
%%%%%%%%%% From Li paper %%%%%%%%%%%%%%
%%%%%%%%%%%%%%%%%%%%%%%%%%%%%%%%%%%%%%%%%%%%%



%%%%%%%%%%%%%%%%%%%%%%%%%%%%%%%%%%%%%%%%%%%%%
%%%%%%%%%% From Myers paper %%%%%%%%%%%%%%
%%%%%%%%%%%%%%%%%%%%%%%%%%%%%%%%%%%%%%%%%%%%%



\usepackage{tikz-qtree}
\usepackage{tikz-qtree-compat}


%Package that makes glossing slightly easier.
\RequirePackage{leipzig}
%\RequirePackage{mathtools} % for \mathrlap

% % % Shortcuts (various borrowed from Jason Zentz)
\RequirePackage{xspace}
\xspaceaddexceptions{]\}}
\usepackage{langsci-textipa}
\usepackage{langsci-branding}
\usepackage{tabto}

%    %%%%%%%%%%%%%%%%%%%%%%%%%%%%%%%%%%%%%%%%%%%%%
%%%%%%%%%% From LSP %%%%%%%%%%%%%%%%%%%%%%%%%
%%%%%%%%%%%%%%%%%%%%%%%%%%%%%%%%%%%%%%%%%%%%%


% \newcommand*{\orcid}{}


\makeatletter
\let\thetitle\@title
\let\theauthor\@author
\makeatother

\newcommand{\togglepaper}[1][0]{
   \bibliography{../localbibliography}
   \papernote{\scriptsize\normalfont
     \theauthor.
     \titleTemp.
     To appear in:
     E. Di Tor \& Herr Rausgeberin (ed.).
     Booktitle in localcommands.tex.
     Berlin: Language Science Press. [preliminary page numbering]
   }
   \pagenumbering{roman}
   \setcounter{chapter}{#1}
   \addtocounter{chapter}{-1}
}

\newbool{bookcompile}
\booltrue{bookcompile}
\newcommand{\bookorchapter}[2]{\ifbool{bookcompile}{#1}{#2}}


%%%%%%%%%%%%%%%%%%%%%%%%%%%%%%%%%%%%%%%%%%%%%
%%%%%%%%%% From ACAL LaTeX template %%%%%%%%%
%%%%%%%%%%%%%%%%%%%%%%%%%%%%%%%%%%%%%%%%%%%%%


\usepackage{tikz}

\newcommand{\circled}[1]{\begin{tikzpicture}[baseline=(word.base)]
\node[draw, rounded corners, text height=8pt, text depth=2pt, inner sep=2pt, outer sep=0pt, use as bounding box] (word) {#1};
\end{tikzpicture}
}


%%%%%%%%%%%%%%%%%%%%%%%%%%%%%%%%%%%%%%%%%%%%%%
%%%%%%%%%%%%%%%%%%%%%%%%%%%%%%%%%%%%%%%%%%%%%%

% Useful Ling Shortcuts


%This makes the \emptyset command be a nicer one
\let\oldemptyset\emptyset
\let\emptyset\varnothing
\newcommand{\nothing}{$\emptyset$}

%Not all of these are explained in the Quick Reference Guide, but they are here bc they are relevant to some of our students. 
\newcommand{\xbar}{\ensuremath{'}\xspace}
\newcommand{\xhead}{\ensuremath{^\circ}\xspace}
\newcommand{\Lb}[1]{$\text{[}_{\text{#1}}$ } %A more convenient left bracket
\newcommand{\Rb}[1]{$\text{]}_{\text{#1}}$ } %A more convenient left bracket
\newcommand{\gap}{\underline{\hspace{1.2em}}}
\newcommand{\vP}{\emph{v}P}
\newcommand{\lilv}{\emph{v}}
\newcommand{\Abar}{A$'$-} %A more convenient A-bar notation
\newcommand{\ph}{$\varphi$\xspace} %A more convenient phi
\newcommand{\pro}{\emph{pro}\xspace}
\newcommand{\subs}[1]{\textsubscript{#1}} %A more convenient subscript
\newcommand{\spells}{$\Longleftrightarrow$} %spellout arrow for morph spellout rules
\newcommand{\tr}[1]{\textit{t}\textsubscript{\textit{#1}}} %easy traces with subscript
\newcommand{\supers}[1]{\textsuperscript{#1}}

%These are borrowed/modified from Andrew Mackenzie's ling-macros package. We have copied them in here (instead of just using the package itself) to avoid a minor clash that was generating an error.


% Abbreviations for glossing, based on Leipzig
\newleipzig{hab}{hab}{habitual}
\newleipzig{rem}{rem}{remote}
\newleipzig{sm}{sm}{subject marker}
\newleipzig{t}{t}{tense}
\newleipzig{aa}{aa}{anti-agreement}
\newleipzig{pron}{pron}{pronoun}
\newleipzig{rec}{rec}{recent}
\newleipzig{om}{om}{object marker}
%\newleipzig{ipfv}{ipfv}{imperfective}
\newleipzig{asp}{asp}{aspect}
\newleipzig{lk}{lk}{linker}
\newleipzig{pcl}{pcl}{particle}
\newleipzig{stat}{stat}{stative}
\newleipzig{ints}{ints}{intensive}
\newleipzig{ascl}{ascl}{assertive subject clitic}
\newleipzig{nascl}{nascl}{non-assertive subject clitic}
\newleipzig{ta}{ta}{tense and/or aspect}
\newleipzig{assoc}{assoc}{associative marker}
\newleipzig{hon}{hon}{honorific}
%\newleipzig{whprt}{wh}{\wh particle}
\newleipzig{sa}{sa}{subject agreement}
\newleipzig{conj}{conj}{conjunction}
%\newleipzig{loc}{loc}{locative}
\newleipzig{expl}{expl}{expletive}
\newleipzig{rcm}{rcm}{reciprocal marker}
\newleipzig{pers}{pers}{persistive}
\newleipzig{fv}{fv}{final vowel}
%\newleipzig{}{}{} %this is just to copy for when I want to add more


%%%%%%%%%%%%%%%%%%%%%%%%%%%%%%%%%%%%%%%%%%%%%%%%%%%%%
%%%%%%%%%%%%%%%%%%%%%%%%%%%%%%%%%%%%%%%%%%%%%%%%%%%%%


%%%%%%%%%%%%%%%%%%%%%%%%%%%%%%%%%%%%%%%%%%%%%
%%%%%%%%%% From Gluckman paper %%%%%%%%%%%%%%
%%%%%%%%%%%%%%%%%%%%%%%%%%%%%%%%%%%%%%%%%%%%%

\newcommand{\pcite}[1]{\citeauthor{#1}'s (\citeyear{#1})}


%%%%%%%%%%%%%%%%%%%%%%%%%%%%%%%%%%%%%%%%%%%%%
%%%%%%%%%% From Myers paper %%%%%%%%%%%%%%
%%%%%%%%%%%%%%%%%%%%%%%%%%%%%%%%%%%%%%%%%%%%%

%% hspace over width of something without showing it
\newcommand*{\hspaceThis}[1]{\settowidth{\LSPTmp}{#1}\hspace*{\LSPTmp}}
% use \hphantom{} for this


%%%%%%%%%%%%%%%%%%%%%%%%%%%%%%%%%%%%%%%%%%%%%
%%%%%%%%%% From Matte paper %%%%%%%%%%%%%%%%%
%%%%%%%%%%%%%%%%%%%%%%%%%%%%%%%%%%%%%%%%%%%%%

%mjkd commented this out in case it's breaking something right now. 
% \newcounter{xlistex}
% \newcommand{\footnoteex}{\stepcounter{xlistex}(\roman{xlistex})}

\newleipzig{sp}{sp}{speaker} %a new leipzig glossing command


%%%%%%%%%%%%%%%%%%%%%%%%%%%%%%%%%%%%%%%%%%%%%
%%%%%%%%%% From GamFranich paper %%%%%%%%%%%%%%
%%%%%%%%%%%%%%%%%%%%%%%%%%%%%%%%%%%%%%%%%%%%%

%MJKD commented these out - both will be provided by the overall LSP book template 

% \bibliography{localbibliography}
% \newcommand{\orcid}[1]{}

%%%%%%%%%%%%%%%%%%%%%%%%%%%%%%%%%%%%%%%%%%%%%
%%%%%%%%%% From Diercks paper %%%%%%%%%%%%%%
%%%%%%%%%%%%%%%%%%%%%%%%%%%%%%%%%%%%%%%%%%%%%

\newcommand{\abar}{$\bar{\text{A}}$\xspace} % this one is different to the \Abar command in Mike's shortcuts doc
\newcommand{\topFeat}{[\textsc{top}]\xspace}


%%%%%%%%%%%%%%%%%%%%%%%%%%%%%%%%%%%%%%%%%%%%%
%%%%%%%%%% From Kinchsular paper %%%%%%%%%%%%%%
%%%%%%%%%%%%%%%%%%%%%%%%%%%%%%%%%%%%%%%%%%%%%

%%%%%%%%%%%%%%%%%%%%%%%%%%%%%%%%%%%%%%%%%%%%%
%%%%%%%%%% From Chen paper %%%%%%%%%%%%%%
%%%%%%%%%%%%%%%%%%%%%%%%%%%%%%%%%%%%%%%%%%%%%

\newcounter{savedex}
\makeatletter
\newcommand*{\saveex}{\setcounter{savedex}{\value{\@xnumctr}}}
\newcommand*{\resumeex}{\setcounter{\@xnumctr}{\value{savedex}}}
\makeatother


%    \input{../localhyphenation}
%     \bibliography{localbibliography}
%     \togglepaper[23]
% }{}

% %custom footer for preprints
% \papernote{\scriptsize\normalfont
%     Change author.
%     Change title.
%     To appear in:
%     Change Volume Editor.  
%     Change volume title.  
%     Berlin: Language Science Press. [preliminary page numbering]
% }



%%%%%%%%%%%%%%%%%%%%%%%%%%%%%%%%%%%%%%%%%%%%%
%%%%%%%%%% From Feibig paper %%%%%%%%%%%%%%
%%%%%%%%%%%%%%%%%%%%%%%%%%%%%%%%%%%%%%%%%%%%%

\usepackage{tikz-qtree}
\usepackage{tikz-qtree-compat}

%%%%%%%%%%%%%%%%%%%%%%%%%%%%%%%%%%%%%%%%%%%%%
%%%%%%%%%% From Olson paper %%%%%%%%%%%%%%
%%%%%%%%%%%%%%%%%%%%%%%%%%%%%%%%%%%%%%%%%%%%%

%MJKD - TIPA is forbidden lol. Commenting out for now but we'll have to address this in the Olson paper

%\usepackage[tone]{tipa}

%%%%%%%%%%%%%%%%%%%%%%%%%%%%%%%%%%%%%%%%%%%%%
%%%%%%%%%% From Caragine paper %%%%%%%%%%%%%%
%%%%%%%%%%%%%%%%%%%%%%%%%%%%%%%%%%%%%%%%%%%%%

\usepackage{placeins}
% \usepackage{graphicx}
% \usepackage{float} %Overleaf suggested this as a solution


%%%%%%%%%%%%%%%%%%%%%%%%%%%%%%%%%%%%%%%%%%%%%
%%%%%%%%%% From Kieffer paper %%%%%%%%%%%%%%
%%%%%%%%%%%%%%%%%%%%%%%%%%%%%%%%%%%%%%%%%%%%%

\usepackage{array}

%%%%%%%%%%%%%%%%%%%%%%%%%%%%%%%%%%%%%%%%%%%%%
%%%%%%%%%% From Payne paper %%%%%%%%%%%%%%
%%%%%%%%%%%%%%%%%%%%%%%%%%%%%%%%%%%%%%%%%%%%%

\usepackage{enumerate}
\usepackage{array,ragged2e}
\usepackage{pgfplots}


%%%%%%%%%%%%%%%%%%%%%%%%%%%%%%%%%%%%%%%%%%%%%
%%%%%%%%%% From Diercks paper %%%%%%%%%%%%%%
%%%%%%%%%%%%%%%%%%%%%%%%%%%%%%%%%%%%%%%%%%%%%

% \usepackage{cgloss}



%%%%%%%%%%%%%%%%%%%%%%%%%%%%%%%%%%%%%%%%%%%%%
%%%%%%%%%% From Li paper %%%%%%%%%%%%%%
%%%%%%%%%%%%%%%%%%%%%%%%%%%%%%%%%%%%%%%%%%%%%



%%%%%%%%%%%%%%%%%%%%%%%%%%%%%%%%%%%%%%%%%%%%%
%%%%%%%%%% From Myers paper %%%%%%%%%%%%%%
%%%%%%%%%%%%%%%%%%%%%%%%%%%%%%%%%%%%%%%%%%%%%



\usepackage{tikz-qtree}
\usepackage{tikz-qtree-compat}


%Package that makes glossing slightly easier.
\RequirePackage{leipzig}
%\RequirePackage{mathtools} % for \mathrlap

% % % Shortcuts (various borrowed from Jason Zentz)
\RequirePackage{xspace}
\xspaceaddexceptions{]\}}
\usepackage{langsci-textipa}
\usepackage{langsci-branding}
\usepackage{tabto}

   %%%%%%%%%%%%%%%%%%%%%%%%%%%%%%%%%%%%%%%%%%%%%
%%%%%%%%%% From LSP %%%%%%%%%%%%%%%%%%%%%%%%%
%%%%%%%%%%%%%%%%%%%%%%%%%%%%%%%%%%%%%%%%%%%%%


% \newcommand*{\orcid}{}


\makeatletter
\let\thetitle\@title
\let\theauthor\@author
\makeatother

\newcommand{\togglepaper}[1][0]{
   \bibliography{../localbibliography}
   \papernote{\scriptsize\normalfont
     \theauthor.
     \titleTemp.
     To appear in:
     E. Di Tor \& Herr Rausgeberin (ed.).
     Booktitle in localcommands.tex.
     Berlin: Language Science Press. [preliminary page numbering]
   }
   \pagenumbering{roman}
   \setcounter{chapter}{#1}
   \addtocounter{chapter}{-1}
}

\newbool{bookcompile}
\booltrue{bookcompile}
\newcommand{\bookorchapter}[2]{\ifbool{bookcompile}{#1}{#2}}


%%%%%%%%%%%%%%%%%%%%%%%%%%%%%%%%%%%%%%%%%%%%%
%%%%%%%%%% From ACAL LaTeX template %%%%%%%%%
%%%%%%%%%%%%%%%%%%%%%%%%%%%%%%%%%%%%%%%%%%%%%


\usepackage{tikz}

\newcommand{\circled}[1]{\begin{tikzpicture}[baseline=(word.base)]
\node[draw, rounded corners, text height=8pt, text depth=2pt, inner sep=2pt, outer sep=0pt, use as bounding box] (word) {#1};
\end{tikzpicture}
}


%%%%%%%%%%%%%%%%%%%%%%%%%%%%%%%%%%%%%%%%%%%%%%
%%%%%%%%%%%%%%%%%%%%%%%%%%%%%%%%%%%%%%%%%%%%%%

% Useful Ling Shortcuts


%This makes the \emptyset command be a nicer one
\let\oldemptyset\emptyset
\let\emptyset\varnothing
\newcommand{\nothing}{$\emptyset$}

%Not all of these are explained in the Quick Reference Guide, but they are here bc they are relevant to some of our students. 
\newcommand{\xbar}{\ensuremath{'}\xspace}
\newcommand{\xhead}{\ensuremath{^\circ}\xspace}
\newcommand{\Lb}[1]{$\text{[}_{\text{#1}}$ } %A more convenient left bracket
\newcommand{\Rb}[1]{$\text{]}_{\text{#1}}$ } %A more convenient left bracket
\newcommand{\gap}{\underline{\hspace{1.2em}}}
\newcommand{\vP}{\emph{v}P}
\newcommand{\lilv}{\emph{v}}
\newcommand{\Abar}{A$'$-} %A more convenient A-bar notation
\newcommand{\ph}{$\varphi$\xspace} %A more convenient phi
\newcommand{\pro}{\emph{pro}\xspace}
\newcommand{\subs}[1]{\textsubscript{#1}} %A more convenient subscript
\newcommand{\spells}{$\Longleftrightarrow$} %spellout arrow for morph spellout rules
\newcommand{\tr}[1]{\textit{t}\textsubscript{\textit{#1}}} %easy traces with subscript
\newcommand{\supers}[1]{\textsuperscript{#1}}

%These are borrowed/modified from Andrew Mackenzie's ling-macros package. We have copied them in here (instead of just using the package itself) to avoid a minor clash that was generating an error.


% Abbreviations for glossing, based on Leipzig
\newleipzig{hab}{hab}{habitual}
\newleipzig{rem}{rem}{remote}
\newleipzig{sm}{sm}{subject marker}
\newleipzig{t}{t}{tense}
\newleipzig{aa}{aa}{anti-agreement}
\newleipzig{pron}{pron}{pronoun}
\newleipzig{rec}{rec}{recent}
\newleipzig{om}{om}{object marker}
%\newleipzig{ipfv}{ipfv}{imperfective}
\newleipzig{asp}{asp}{aspect}
\newleipzig{lk}{lk}{linker}
\newleipzig{pcl}{pcl}{particle}
\newleipzig{stat}{stat}{stative}
\newleipzig{ints}{ints}{intensive}
\newleipzig{ascl}{ascl}{assertive subject clitic}
\newleipzig{nascl}{nascl}{non-assertive subject clitic}
\newleipzig{ta}{ta}{tense and/or aspect}
\newleipzig{assoc}{assoc}{associative marker}
\newleipzig{hon}{hon}{honorific}
%\newleipzig{whprt}{wh}{\wh particle}
\newleipzig{sa}{sa}{subject agreement}
\newleipzig{conj}{conj}{conjunction}
%\newleipzig{loc}{loc}{locative}
\newleipzig{expl}{expl}{expletive}
\newleipzig{rcm}{rcm}{reciprocal marker}
\newleipzig{pers}{pers}{persistive}
\newleipzig{fv}{fv}{final vowel}
%\newleipzig{}{}{} %this is just to copy for when I want to add more


%%%%%%%%%%%%%%%%%%%%%%%%%%%%%%%%%%%%%%%%%%%%%%%%%%%%%
%%%%%%%%%%%%%%%%%%%%%%%%%%%%%%%%%%%%%%%%%%%%%%%%%%%%%


%%%%%%%%%%%%%%%%%%%%%%%%%%%%%%%%%%%%%%%%%%%%%
%%%%%%%%%% From Gluckman paper %%%%%%%%%%%%%%
%%%%%%%%%%%%%%%%%%%%%%%%%%%%%%%%%%%%%%%%%%%%%

\newcommand{\pcite}[1]{\citeauthor{#1}'s (\citeyear{#1})}


%%%%%%%%%%%%%%%%%%%%%%%%%%%%%%%%%%%%%%%%%%%%%
%%%%%%%%%% From Myers paper %%%%%%%%%%%%%%
%%%%%%%%%%%%%%%%%%%%%%%%%%%%%%%%%%%%%%%%%%%%%

%% hspace over width of something without showing it
\newcommand*{\hspaceThis}[1]{\settowidth{\LSPTmp}{#1}\hspace*{\LSPTmp}}
% use \hphantom{} for this


%%%%%%%%%%%%%%%%%%%%%%%%%%%%%%%%%%%%%%%%%%%%%
%%%%%%%%%% From Matte paper %%%%%%%%%%%%%%%%%
%%%%%%%%%%%%%%%%%%%%%%%%%%%%%%%%%%%%%%%%%%%%%

%mjkd commented this out in case it's breaking something right now. 
% \newcounter{xlistex}
% \newcommand{\footnoteex}{\stepcounter{xlistex}(\roman{xlistex})}

\newleipzig{sp}{sp}{speaker} %a new leipzig glossing command


%%%%%%%%%%%%%%%%%%%%%%%%%%%%%%%%%%%%%%%%%%%%%
%%%%%%%%%% From GamFranich paper %%%%%%%%%%%%%%
%%%%%%%%%%%%%%%%%%%%%%%%%%%%%%%%%%%%%%%%%%%%%

%MJKD commented these out - both will be provided by the overall LSP book template 

% \bibliography{localbibliography}
% \newcommand{\orcid}[1]{}

%%%%%%%%%%%%%%%%%%%%%%%%%%%%%%%%%%%%%%%%%%%%%
%%%%%%%%%% From Diercks paper %%%%%%%%%%%%%%
%%%%%%%%%%%%%%%%%%%%%%%%%%%%%%%%%%%%%%%%%%%%%

\newcommand{\abar}{$\bar{\text{A}}$\xspace} % this one is different to the \Abar command in Mike's shortcuts doc
\newcommand{\topFeat}{[\textsc{top}]\xspace}


%%%%%%%%%%%%%%%%%%%%%%%%%%%%%%%%%%%%%%%%%%%%%
%%%%%%%%%% From Kinchsular paper %%%%%%%%%%%%%%
%%%%%%%%%%%%%%%%%%%%%%%%%%%%%%%%%%%%%%%%%%%%%

%%%%%%%%%%%%%%%%%%%%%%%%%%%%%%%%%%%%%%%%%%%%%
%%%%%%%%%% From Chen paper %%%%%%%%%%%%%%
%%%%%%%%%%%%%%%%%%%%%%%%%%%%%%%%%%%%%%%%%%%%%

\newcounter{savedex}
\makeatletter
\newcommand*{\saveex}{\setcounter{savedex}{\value{\@xnumctr}}}
\newcommand*{\resumeex}{\setcounter{\@xnumctr}{\value{savedex}}}
\makeatother


   \input{../localhyphenation}
   \boolfalse{bookcompile}
   \togglepaper[23]%%chapternumber
}{}

\begin{document}
\maketitle

\section{Introduction}\label{sec:feibig:sec1}

Kalenjin is a rather understudied language cluster. First grammar descriptions of the relevant varieties can be found in \citet{Toweet79}, \citet{Rottland82}, \citet{Crazzolara} and \citet{Beroja}. Theoretical papers on Kipsigis include \citet{Kouneli2020} on the tripartite number marking system in Kipsigis, \citet{Kouneli2022} on the so called inflection classes in Kipsigis, and \citet{Bossi23} on epistemic modality. This paper provides the very first description of verbal number in two Kalenjin varieties, Kipsigis and Pokot. By doing so, this paper shows how Kalenjin verbal number is different from how verbal number works in many other languages. While many languages mark both event and participant number by the very same marking strategy, Kalenjin exhibits two distinct marking strategies for the two grammatical concepts.\\ This paper is structured as follows: In section two, first, the topic of pluractionality is introduced briefly. The second section further introduces the languages and what pluractionality looks like in each of them. The variety of semantic functions of Kalenjin pluractionality is demonstrated via many examples from Pokot.  
The third section provides a morphological description of Kipsigis and Pokot pluractionality. The fourth section introduces the distinction of event and participant number and argues against the conflation of event and participant number in Kipsigis and Pokot and why this is problematic for existing generative accounts of verbal number. The last section concludes. 

\section{Description of pluractionality in Kipsigis and Pokot}\label{sec:feibig:sec2}

While nominal number, an obviously related category, marks the number of individuals, verbal number marks the number of events (or situations, in the words of \citealt{Mattiola2019}). Verbal number can be defined as in \REF{ex:fiebig:definition}. 

\ea \label{ex:fiebig:definition} \textbf{Verbal number}\\
``A phenomenon that marks plurality or multiplicity of the situations (i.e. states and events) encoded by the verb through any morphological mean that modifies the form of the verb itself.'' \citep[4]{Mattiola2019}
\z


\noindent Following \citet{Corbett2000}, verbal number can be observed in two forms: participant number (multiple participants create a multiplicity of events) and event number (a single participant creates/experiences a multiplicity of events). This distinction will be discussed in section \sectref{sec:fiebig:partNum} in detail. I follow \cite{Henderson2012} and others in calling the second flavour (event number) \textit{pluractionality} and also adapt his notation of calling the morphemes which derive this multiple event readings \textit{pluractionals}.\\
Verbal number is most common in North America but can also be observed in several language families in South America, the Pacific region, Caucasus and all four major families of Africa \citep{wals-80}. In Africa, it can be observed most strongly in the Nilo-Saharan language family, which is the language family Kalenjin also belongs to. Kipsigis and Pokot are two varieties of Kalenjin, a dialect cluster belonging to the Southern Nilotic branch of Nilo-Saharan. 
Kipsigis is spoken by approximately 2 million speakers and Pokot by approximately 778,000 speakers in Kenya \citep{Ethnologue2023}.
Although the languages are closely related they are not mutually intelligible.
If not otherwise indicated, all data from Kipsigis and Pokot in this paper come from fieldwork and were obtained in elicitations with five speakers in Kenya.


\subsection{Kipsigis}


Kipsigis verb stems can be reduplicated to derive a pluractional meaning. Pluractionality in Kipsigis is productive and can be used with Kipsigis words and also Swahili loan words. It doesn't seem to be limited to certain verb classes, but no Aktionsart-tests exist for Kipsigis yet, so a limitation to certain Aktionsart-classes can't be excluded at this point. In my field notes, only two verbs were unable to be reduplicated: 
\textit{surie} `to hate' and \textit{wɛndat} `to walk'.
The meaning of Kipsigis reduplication differs slightly from root to root. Some of the meanings are `to do something in stages', `do something again and again', or `do something on and off'. The core meaning all verbal reduplications share seems to be `this event takes place more than once'. See \REF{ex:fiebig:firstKip} for a first Kipsigis example.


\ea Kipsigis\footnote{Tone is indicated whenever possible, with exception of a few examples which were obtained in online elicitations where transcription was difficult because of unstable audio quality.} \label{ex:fiebig:firstKip}
\begin{xlist}

\ex \label{ex:fiebig:firstKipA}
\gll í-ít-è\\
\textsc{cl2}-count-\textsc{ipfv}\\
\glt `he counts'

\ex \label{ex:fiebig:firstKipB}
\gll í-ít$\sim$ɑ̀ː$\sim$ìt-ì \\
\textsc{cl2}-count$\sim$\textsc{lk}$\sim$count-\textsc{ipfv}\\
\glt `He is counting several times.'\\ Comment: ``You count and when you finish you start counting again. Maybe to confirm the sum you counted the first time.''

\end{xlist}

\z 


\noindent In \REF{ex:fiebig:firstKipB} the verbal stem \textit{ít} `count' is reduplicated. In between the two reduplicative elements, a linking vowel \textit{ɑː} appears, as is the case regularly with most short-vowel-stems in reduplication.
The meaning of the reduplicated verb is given in the translation of \REF{ex:fiebig:firstKipB} and the speaker comment indicates that the whole event is repeated in this case and not only subparts of it. As shown in \REF{ex:fiebig:first}, Kipsigis' reduplications are incompatible with modifiers like \textit{kɔnɪl} \textsl{a}\textit{ɣɛːŋɛ} `once'.

\ea \label{ex:fiebig:first}
\gll pir∼ɑː∼pir-e laːkwɛːt kibɛːt (*kɔnɪl aɣɛːŋɛ)\\
beat∼\textsc{lk}∼beat-\textsc{ipfv} child Kibeet times one\\
\glt `Kibeet is beating the child again and again (*once).'\\ 
Comment: ``That does not work, for pir-aː-pire you don't know how often it happens, but not once.''
\z

\noindent This incompatibility, in addition to the example's translations and the speaker's comment, illustrates how the meaning of pluractionality is `the event takes place more than once'.





\subsection{Pokot}
\subsubsection{Basic facts of Pokot pluractionality}
Unlike Kipsigis, Pokot has two pluractional morphemes: reduplication and \textit{-sete}\footnote{I am following \cite{Marantz82} here in assuming reduplication to be a morpheme.}. The morpheme \textit{-sete} sometimes appears in slightly varying vowel quality due to phonological processes. 
Both morphemes can appear on the same stem, in both possible orders (verb-\textsc{red}-\textit{sete} or verb-\textit{sete}-\textsc{red}). The orders are abstractly illustrated in \REF{ex:fiebig:abstr}. For the same verb, pluractional expressions differ in meaning between the morphemes and sometimes between the possible combinations of them. 

\ea co-occuring pluractional morphemes \label{ex:fiebig:abstr}
\begin{xlist}
\ex verb-\textsc{red}-\textit{sete}: verb$\sim$verb-\textit{sete} \label{ex:fiebig:abstra}
\ex verb-\textit{sete}-\textsc{red}: verb-\textit{sete}$\sim$verb-\textit{sete} 
\end{xlist}
\z

Both morphemes share a meaning like `this event happens more than once', but the exact meaning of each of the morphemes and each of the combinations differs with the specific verb. See \sectref{sec:fiebig:morphKal} for concrete examples of this variable ordering.


\ea Pokot \label{ex:fiebig:telicity}
\begin{xlist}
\ex 
\gll ì-wù:t-éɲ-ì\\
\textsc{2sg}-shoot-\textsc{2sg-ipfv}\\
\glt `You will shoot' or `you are shooting.'
\ex 
\gll ì:-wú:t$\sim$wú:t-éːɲ-ì\\
\textsc{2sg}-\textbf{shoot$\sim$shoot}-\textsc{2sg-ipfv}\\
\glt `You are just shooting and shooting without aiming.' \label{ex:fiebig:-red}
\ex 
\gll í-wùt-sə́tə̀-nɛ́ːɲ-ì\\
2\textsc{sg}-shoot-\textbf{sete}-2\textsc{sg-ipfv}\\
\glt `You are shooting so many times.' / `It is your habit to shoot a lot.' \label{ex:fiebig:-sete}
\end{xlist}
\z

\noindent Example \REF{ex:fiebig:-red} shows a reduplication of the verbal stem \textit{wù:t}. The translation of \REF{ex:fiebig:-red} indicates a repeated event of shooting. Modifying the verbal stem with the pluractional suffix \textit{-sete} in \REF{ex:fiebig:-sete} also derives meanings where the event takes place several times. Both seem to differ in their interpretation only slightly and sometimes not at all, depending on the specific verb. 
Often they seem to differ in intensity, as in \REF{ex:fiebig:supplPok}, in the specific amount of repetitions as in \REF{ex:fiebig:varyB}/\REF{ex:fiebig:vary} or even in telicity, as in \REF{ex:fiebig:telicity}. 
Both Pokot pluractionals, just as reduplication in Kipsigis, generally cannot co-occur with modifiers conveying the meaning `once' (although, see an exceptional reading in example \REF{ex:fiebig:freq}). As \REF{ex:fiebig:PokOnce} shows, neither reduplication nor \textit{-sete} is compatible with \textit{kəɲəl akoŋge} `once.'


\ea Pokot \label{ex:fiebig:PokOnce}

    \begin{xlist}
        \ex 
        \gll kì-pìs$\sim$ɑ̀ː$\sim$pís ɲíndə̀ (*kəɲəl akoŋge) \\
        \textsc{pst}-split$\sim$\textsc{lk}$\sim$split \textsc{pron}.\textsc{3sg} times one \\
        \glt `He split again and again (*once).' \\
        Comment: ``It can't be once because it is repetition.''
        \ex 
        \gll kì-pìs-sàtə́ kwén (*kəɲəl akoŋge) \\
        \textsc{pst}-split-sete firewood times one \\
        \glt `He split firewood again and again (*once).' \\
    \end{xlist}

\z


\noindent \REF{ex:fiebig:PokOnce} also shows that the same linking vowel as in Kipsigis (\textit{-aː-}) also appears with short-vowel-stems in Pokot.\\ 
Just as in Kipsigis, pluractionality in Pokot is productive and can be used with Pokot and also Swahili loan words. It doesn't seem to be limited to certain verb classes, but no Aktionsart-tests exist for Pokot yet, so a limitation to certain Aktionsart-classes can't be excluded at this point. In my field notes, there is no word that couldn't be modified by reduplication and \textit{-sete}.

\subsubsection{Semantic functions of Pokot pluractionality}\label{sec:fiebig:semantics}

The semantics of Kalenjin pluractionality varies from root to root. To illustrate the great variety of semantic functions I matched my data with the cross- linguistically observable semantic functions of pluractionality. \cite{Mattiola2019} provides a typological investigation of pluractionality (246 languages, based on \citealt{wals-80}) and as a result provides a list of semantic functions that can be expressed by pluractionality. Although Kipsigis and Pokot exhibit a variety of semantic functions, for space reasons, only Pokot examples will be discussed in detail in the following to illustrate the semantic variety of Kalenjin pluractionality. \\
Pokot's two pluractional morphemes are able to express at least 5 out of the 17 possible functions \cite{Mattiola2019} describes. Out of the several semantic functions observable cross-linguistically but not in Pokot (and Kipsigis), two functions will be discussed in some detail: Habituality is discussed because several translations seem to indicate this function for Pokot pluractionality, but Pokot actually uses imperfective aspect to express habituality. Participant number, on the other hand, is taken to be one of the two flavours of verbal number (e.g. \citealt{Corbett2000}, \citealt{Amato2018}, \citealt{Thornton2019}) but can be shown to be two distinct categories in Pokot and Kipsigis. Arguments against conflating these categories in Kipsigis are shown in the next section. \\
In order to describe the semantic functions of pluractionality, \cite{Mattiola2019} uses the following notations, based on \citet[77]{Cusic81}, which will appear also in this paper:

\ea \textbf{event notation} \citep[21]{Mattiola2019}

\begin{xlist}

\ex situation: \\ an event or state

\ex phase: \\ the subpart of a situation

\ex occasion: \\ a specific occurrence of a situation: the situation with its specific time frame, place and involved participants

\end{xlist}

\z


\subsubsubsection{Iterativity}

\textit{Iterativity}, one of the core meanings of pluractionality, is described as the sequential repetition of an event on a particular occasion \citep[22]{Mattiola2019}. Both, reduplication and \textit{-sete} can express this meaning in Pokot.  

\ea Pokot
\begin{xlist}
    
\ex 
\gll lùw-sə́tèː-nə́ mɛ́ːsà\\
hit.\textsc{imp}-sete-\textsc{ipfv}? table\\
\glt `hit the table repeatedly/continuously without a reason!'
\ex 
\gll lùw$\sim$á$\sim$lúw-ɔ́ʰ mɛ́ːsà\\
hit$\sim$\textsc{lk}$\sim$hit-\textsc{imp} table\\
\glt `hit the table repeatedly for the sake of something!'\\
Comment: e.g. to follow a rhythm or to a baby to get it to fall asleep.

\end{xlist}
\z

\subsubsubsection{Frequentative}



Following \citet[24]{Mattiola2019}, \textit{frequentative} is the repetition of a specific situation, performed over multiple different occasions. The downtime between the repetitions is long enough to understand each as separate occasions. This matches the so called \textit{plurality of and in events}, as described by \cite{Cusic81}. The translations of \REF{ex:fiebig:freq} show that Pokot reduplication and \textit{-sete} are able to convey this meaning. 

\ea Pokot \label{ex:fiebig:freq}
\begin{xlist}
\ex
\gll kì-tɛ́k$\sim$à$\sim$tɛ̀k-é-j kɔ̀ɲíl ákɔ̀ŋə́\\
\textsc{pst}-build$\sim$\textsc{lk}$\sim$build-\textsc{3sg-ipfv} times one\\
\glt `There were several instances of him building once.'
\ex
\gll kì-tɛ́k-sàtɔ́-j kɔ̀ɲíl ákɔ̀ŋə́\\
\textsc{pst}-build-sete-\textsc{ipfv} times one\\
\glt `There were many instances of him building once.'
\end{xlist}
\z

\noindent The building takes place only once at each occasion. Since this is repeated over a bigger time frame there is still a repetitive meaning, the \textit{Frequentative} reading. This reading gives rise to the licit combination of pluractionality and \textit{kɔ̀ɲíl ákɔ̀ŋə́} `once' which is illicit for other verbs, like \textit{pɪs} `split' in \REF{ex:fiebig:PokOnce}. That a frequentative reading is not possible for all verbs, again shows how the interpretation of pluractionality is highly dependent on the lexical semantics of each root.

\subsubsubsection{Spatial distribution}

\textit{Spatial distribution} describes multiple events which are distributed over several places \citep[24]{Mattiola2019}. \cite{Mattiola2019} doesn't explain if these interpretations are necessary with pluractionality in the languages he discusses or merely a possible situation in which the situation described by pluractionality can occur. For example, if a verb like `to build' is reduplicated, spatial distribution is a necessary interpretation (if the thing being built is not built, torn down, and built again etc.). The Pokot example in \REF{ex:fiebig:spatz} shows such an interpretation: The translation of the example with `here and there' indicates a spatial distribution of the event.

\ea Pokot reduplication \label{ex:fiebig:spatz} \\
\gll kàríb̀ʊ̀ kì-jɛ́ɣ$\sim$à$\sim$jɛ́ɣ nà:djə́ sú:mì\\ 
almost \textsc{pst.dist}-drink$\sim$\textsc{lk}$\sim$drink Nadja poison\\
\glt `Nadja almost drank poison here and almost drank poison there.' 
\z 

\noindent But spatial distribution is only one possible context in which the pluractionality can be interpreted. The context of \REF{ex:fiebig:water-glass} specifies a non-spatially-distributed reading and the very same pluractional form \textit{kì:-jɛ́ɣ-ɑ̀-jɛ́ɣ} `(s)he drank again and again' is felicitous in this context as well. This means that iterative or frequentative meanings are always conveyed by pluractionality, but spatial distribution is only a possible context of this repetition. 

\ea \label{ex:fiebig:water-glass} Context: There is a person having several glasses of water in front of him. He is drinking them in a specific fashion: he drinks each glass really really fast, it only takes him a few seconds. But then he waits a good while to drink the next glass. I ask you to describe this behaviour to me afterwards.
\begin{xlist}
\ex \gll kì:-jɛ́ɣ$\sim$ɑ̀$\sim$jɛ́ɣ lɔ̀kɔ̀dɛ́t-ì nàu píɣ là\textsuperscript{w}ɛ́ʰl\\
\textsc{pst.dist}-drink$\sim$\textsc{lk}$\sim$drink cup-\textsc{pl} of water fast\\
\glt `He drank cups of water quickly.'
\end{xlist}
\z

\noindent Interestingly, no \textit{-sete} form appeared in my data which was translated with a similar construction indicating spatial distribution. Further investigations are needed to determine if the (optional) spatial distribution interpretation is maybe limited to reduplication.

\subsubsubsection{Habituality}

\citet[24]{Mattiola2019}, following \citet{Comrie76}, defines \textit{habituality} as a repeated action which is typical of a specific period of time. Several translations of pluractional sentences in Pokot suggest that Pokot pluractionality expresses this semantic function. See for example the a. sentences in \REF{ex:fiebig:builder1} and \REF{ex:fiebig:builder2}, which have a clearly habitual translation. 


\ea Context: We are talking about somebody who was a professional builder. \label{ex:fiebig:builder1}
\begin{xlist}
\ex 
\gll ki-tɛk-satɑ-j ɲində kɔ-rɪn\\
\textsc{pst}-build-sete-\textsc{ipfv} pron.\textsc{3sg} house-\textsc{pl}\\
\glt `He used to build houses.'
\ex 
\gll \#kì-tɛ́k-sàtà ɲində kɔ-rin\\
\textsc{pst}-build-sete pron.\textsc{3sg} house-\textsc{pl}\\
\glt `He built several times.'
\end{xlist}
\z

\ea  Context: Somebody had a career as professional builder a while ago\label{ex:fiebig:builder2}
\begin{xlist}
\ex 
\gll kì-ték$\sim$ɑ̀$\sim$tèk-é-j\\
\textsc{pst}-build$\sim$\textsc{lk}$\sim$build-3-\textsc{ipfv}\\
\glt `He used to build a lot.'
\ex 
\gll \#kì-tɛ́k$\sim$à$\sim$tɛ̀k\\
\textsc{pst}-build$\sim$\textsc{lk}$\sim$build\\ 
\glt `He built many times.'
\end{xlist}
\z
    
\noindent But neither reduplication, nor \textit{-sete} can introduce habituality on its own. Overtly marked imperfectivity is needed additionally to express habituality, as the b. examples of \REF{ex:fiebig:builder1} and \REF{ex:fiebig:builder2} show respectively: Kipsigis and Pokot both mark imperfectivity, while unmarked verbs are perfective (\citealt{Kouneli2020}, \citealt{Toweet79}). The aspectually unmarked, and therefore perfective, versions of the same pluractional constructions are infelicitous in a clearly habitual context. 
Further, habituality can even be introduced by imperfectivity on its own, not only in combination with reduplication or \textit{-sete}: 

\ea
\begin{xlist}
\ex 
\gll kì-rɑ́l-ə́-j\\
\textsc{pst}-cough-\textsc{3-ipfv}\\
\glt `He used to cough.'\\
Comment: ``He had this habit but got better''
\ex
\gll ki-ral\\
\textsc{pst}-cough\\
\glt `He coughed.' (once)
\end{xlist}
\z

\noindent Habituality is, therefore, rather a property of imperfectivity, and not of pluractionality in Pokot.

\subsubsubsection{Event-internal plurality}

\textit{Event-internal plurality} is defined as a single situation consisting of hard to distinguish sub-parts, e.g. ``He whistles.'' (\citealt{Mattiola2019}, see also \citealt{Cusic81}). The pluractional form of \textit{támbìl} `to swim' in \REF{ex:fiebig:swim} provides such an example. As the context shows, not the whole act of swimming from one point to another is repeated. Instead, and this is what the speaker's comment also implies, the indistinguishable subparts of swimming are repeated. 



\ea Context: The person we are talking about is swimming from one side to the other without taking a rest.\label{ex:fiebig:swim}
\begin{xlist}

\ex
\gll tɑ́mbìl$\sim$ɑ̀$\sim$tɑ̀mbìl-éj pír\\
swim$\sim$\textsc{lk}$\sim$swim-\textsc{ipfv} water\\
\glt `He swims repeatedly.'\\
Comment: ``Swimming means to do the movement with your hands over and over again.''
\ex
\gll  kì-támbìl-sètɔ́\\
\textsc{pst}-swim-sete\\
\glt `He swims repeatedly.'

\end{xlist}
\z


\subsubsubsection{Intensity}

The semantic function \textit{intensity} is described as ``a situation done with more effort or whose result is augmented with respect to the normal happening of the same situation'' \citep[36]{Mattiola2019}. The context of \REF{ex:fiebig:cough}, as well as the translations of \REF{ex:fiebig:cough} and \REF{ex:fiebig:build} indicate that both reduplication and \textit{-sete} can express \textit{intensity} via pluractionality.

\ea Context: The person we are talking about was really sick. They were coughing only once but that coughing lasted for 2 minutes. \label{ex:fiebig:cough}
\begin{xlist}
\ex 
\gll kì-rál-sátà\\
\textsc{pst}-cough-sete\\
\glt `he coughed really badly'
\end{xlist}
\z

\ea \label{ex:fiebig:build}
\gll kì-tɛ́k$\sim$à$\sim$tɛ̀k kí\\
\textsc{pst}-build$\sim$\textsc{lk}$\sim$build house\\
\glt `he built a house faster than normal.' 
\z

\subsubsubsection{Indefiniteness}


Another function that can be observed cross-linguistically and also expressed by Pokot pluractionals is what \citet{Mattiola2019} calls \textit{indefiniteness}. \citet[40]{Mattiola2019} describes this function as a lack of specific orientation or goal. Other analyses treat this function of pluractionality rather as a change in telicity, see e.g. \citet{Henderson2012}. 
\REF{ex:fiebig:shoot-non telic} shows how the reduplicated verb \textit{wúːt} `shoot' expresses a non-telic reading ('without aiming'), while the form modified by \textit{-sete} in \REF{ex:fiebig:non-tel-shoot-b} doesn't express such a non-telic reading.


\ea Pokot \label{ex:fiebig:shoot-non telic}
\begin{xlist}
\ex 
\gll ìː-wúːt$\sim$wúːt-ɛ́ːɲ-ì\\
\textsc{2sg}-shoot$\sim$shoot-\textsc{2sg-ipfv}\\
\glt `You are just shooting and shooting \textbf{without aiming}.'
\ex \label{ex:fiebig:non-tel-shoot-b}
\gll í-wùt-sə́tə̀-nɛ́ːɲ-ì\\
2\textsc{sg}-shoot-sete-2\textsc{sg-ipfv}\\ 
\glt `You are shooting regularly.'
\end{xlist}
\z

\noindent On the other hand, \REF{ex:fiebig:beat-non telic} shows that \textit{-sete} is able to express non-telicity/ indefiniteness as well. In this example, it is the form modified by \textit{-sete} which expresses a non-telic version of the verb, while the reduplicated form of \textit{lùw} `hit' is translated with `for the sake of something', clearly indicating a goal of the event.


\ea Pokot \label{ex:fiebig:beat-non telic}
\begin{xlist}
\ex 
\gll lùw-sə́tèː-nə́ mɛ́ːsà\\
hit.\textsc{imp}-sete-\textsc{ipfv} table\\
\glt `hit the table repeatedly/continuously \textbf{without a reason}!'
\ex 
\gll lùw$\sim$á$\sim$lúw-ɔ́ʰ mɛ́ːsà\\
hit$\sim$\textsc{lk}$\sim$hit-\textsc{imp} table\\
\glt`Hit the table repeatedly for the sake of something.'\\
Comment: e.g. to follow a rhythm or to for a baby to make it fall asleep by listening to you.
\end{xlist}
\z

\noindent Therefore, both reduplication and \textit{-sete} are able to express this function, but apparently always only one of them, depending on the predicate.

\subsubsection{Summary of the semantic functions}

The previous examples show that Pokot is able to express at least 5 out of the 17 functions of pluractionality \citet{Mattiola2019} identifies. \tabref{tab:fiebig:overview} shows what morpheme is able to express which function and which functions are unattested. 

\begin{table}
\begin{tabular}{l  c  c  c}
\lsptoprule
& \textit{-sete} & \textsc{red}\\
\midrule
iterativity & \Checkmark & \Checkmark\\
frequentative & \Checkmark & \Checkmark\\
spatial distribution & (possible context) & no data \\
habitual & \XSolidBrush & \XSolidBrush \\
event-internal & \Checkmark & \Checkmark \\
intensity & \Checkmark & \Checkmark\\
indefiniteness (change in telicity) & \Checkmark & \Checkmark \\
successive events & no data & no data \\
completeness & no data & no data \\
singulactionality & no data & no data \\
emphasis & no data & no data\\
reciprocity & no data & no data\\
generic & no data & no data\\
continuative & no data & no data\\
antipassive & no data & no data \\
causative & no data & no data\\
participant plurality & no data & no data \\
\lspbottomrule
\end{tabular}
\caption{Attested semantic functions of pluractional morphemes in Pokot}\label{tab:fiebig:overview}
\end{table}

\noindent For all other functions, no examples came up in my field notes. This doesn't indicate that these functions don't exist in Pokot, but at least hints at the conclusion that these functions are not as readily available as the attested ones. \textit{Spatial distribution} doesn't seem to be a function Pokot pluractionality expresses, but rather is one of several possible contexts in which repetition can be interpreted. Participant plurality, and how it is not expressed by reduplication or \textit{-sete} is going to be discussed in detail in \sectref{sec:fiebig:partNum}.



\section{Morphological description of Kalenjin pluractionality}\label{sec:fiebig:morphKal}

Cross-linguistically, affixation, reduplication and lexical alternation/suppletion are the most common marking strategies for pluractionality \citet[65]{Mattiola2019}. Pokot and Kipsigis therefore display two of the three most common strategies. Suppletion is also present in Kipsigis verbal number, but only to mark participant number, as will be shown in \sectref{sec:fiebig:partNum}. Why Pokot employs two different marking strategies and how exactly their meaning differs, unfortunately, needs to remain a question for further research at this point.
For both Pokot and Kipsigis, the reduplicative morpheme is structurally located below person, tense and aspect morphology, close to the verbal root but above little \textit{v}, the verbalizing head in theories like distributed morphology \citep{DM}. As \REF{ex:fiebig:PokSize} illustrates, reduplication only targets the verbal stem.  Including person and aspect marking in the reduplication leads to ungrammaticality, as in \REF{ex:fiebig:incorrect}.\footnote{If not indicated otherwise, all following examples are illustrative for the behaviour of Pokot \textit{and} Kipsigis, even if only one language is shown explicitly.}


\ea Pokot  \label{ex:fiebig:PokSize}
\begin{xlist}
\ex
\gll ì-wùːt-éːɲ-ì\\
\textsc{2sg}-shoot-\textbf{\textsc{2sg-ipfv}}\\
\glt `You will shoot' or `you are shooting.'
\ex \label{ex:fiebig:correct}
\gll ìː-wúːt$\sim$wúːt-éːɲ-ì\\
\textsc{2sg}-shoot$\sim$shoot-\textbf{\textsc{2sg-ipfv}}\\
\glt `You are just shooting and shooting without aiming.' 
\ex \label{ex:fiebig:incorrect}
\gll *ì-wùːt-éɲ-ì$\sim$wùːt.éɲ.ì\\
2-shoot-\textbf{2-\textsc{ipfv}}$\sim$shoot.\textsc{ipfv}\\
\glt Intended interpretation: `He is shooting again and again.' 
\end{xlist}
\z

\ea Kipsigis \label{ex:fiebig:KipAsp}
\begin{xlist}
\ex
\gll ɑ́l-é kíbêːt tʃàbàtɪ́\\
buy-\textsc{ipfv} Kibeet chapati\\
\glt`Kibeet is buying chapati.'
\ex \label{ex:fiebig:kipAspA}
\gll *ɑ́l-é$\sim$ɑ́l.é kíbêːt tʃàbàtɪ́\\
buy-\textsc{ipfv}$\sim$buy.\textsc{ipfv} Kibeet chapati\\
\glt`Kibeet is buying chapati over and over again.' 
\ex \label{ex:fiebig:kipAspB} 
\gll ɑ́l$\sim$ɑ̀ːl-é  kíbêːt tʃàbàtɪ́\\
buy$\sim$buy-\textsc{ipfv} Kibeet chapati\\
\glt`Kibeet is buying chapati over and over again.'
\end{xlist}
\z

\noindent The same is true for Kipsigis, as \REF{ex:fiebig:KipAsp} shows. The aspect marker \textit{-e} cannot be reduplicated with the verbal stem (as in \REF{ex:fiebig:kipAspA}) but must be merged after the reduplication as in \REF{ex:fiebig:kipAspB}. Pokot's \textit{-sete} suffix, as well, appears closer to the stem than aspect and person marking as \REF{ex:fiebig:-sete} shows. 
Tense morphology, as well, cannot be reduplicated with the stem. Example \REF{ex:fiebig:tense} shows how doing so leads to ungrammaticality.

\ea Kipsigis \label{ex:fiebig:tense}
\begin{xlist}
\ex
\gll ki-bir$\sim$ɑː$\sim$bir-i\\
\textsc{pst}-beat$\sim$\textsc{lk}$\sim$beat-\textsc{ipfv}\\
\glt `He was beating again and again.'
\ex
\gll *ki-bir-(ɑː-)ki.bir-i\\
\textsc{pst}-beat$\sim$\textsc{lk$\sim$pst}.beat-\textsc{ipfv}\\
\glt `He was beating again and again.'
\end{xlist}
\z


\noindent Other verbal suffixes, as for example applicative morphology, also cannot be reduplicated with the stem as illustrated in \REF{ex:fiebig:KipSize} for Kipsigis. Analogously, \REF{ex:fiebig:PokotBoth} shows how in Pokot applicative morphology applies further away from the root than both reduplication and \textit{-sete}.

\ea Kipsigis \label{ex:fiebig:KipSize} 
\begin{xlist}
\ex
\gll  kɔ̀ː-á-tɛ́m$\sim$áː$\sim$tɛ́m-\fbox{ɛ̀ːn} ìmbàr mògòːmbéːt\\
\textsc{pst2-1sg}-dig$\sim$\textsc{lk}$\sim$\textbf{dig}-\textsc{\textbf{appl}} farm hoe\\
\glt `I dug the farm with a hoe again and again.'
\ex
\gll *kɔ̀ː-á-tɛ́m-ɛːn$\sim$(áː)$\sim$tɛ́m-ɛ̀ːn ìmbàr mògò:mbé:t\\
\textsc{pst2-1sg}-dig-\textsc{\textbf{appl}}$\sim$\textsc{lk}$\sim$\textbf{dig} farm hoe\\
\glt `I dug the farm with a hoe again and again.'
\end{xlist}
\z


\ea Pokot \label{ex:fiebig:PokotBoth} \\
\gll kɛ̀-líl$\sim$á$\sim$líl-sátá-tSíː\\
to-hate$\sim$\textsc{lk}$\sim$\textbf{hate-sete}-\textsc{\textbf{appl}}\\
\glt `to hate someone a lot.'	
\z 

\noindent In contrast to the morphemes in the data above, the verbalizing head \textit{-an} in both Kipsigis and Pokot can be reduplicated and appears closer to the root then \textit{-sete}. 


\ea Pokot \label{ex:fiebig:PokotRain}
\begin{xlist}
    \ex
    \gll rɔ́p\\
	rain\\
	\glt `the rain' (noun)
	\ex
    \gll mak-an rɔ́p-\textbf{an}\\
	want-\textsc{1sg} rain-\textsc{vblz}\\
	\glt `I want it to rain.'
	\ex
    \gll míː rɔ̀p-\textbf{àn}-ə̀-j\\
	\textsc{cop} rain-\textsc{vblz}-3-\textsc{ipfv}\\
	\glt `It is raining.' 
	\ex \label{ex:fiebig:rainRED}
    \gll rɔ́p-\textbf{àn}$\sim$ɑ̀$\sim$\textbf{rɔ̀p.\textbf{àn}}-ə́-j\\
	rain-\textsc{vblz}$\sim$\textsc{lk}$\sim$rain.\textsc{vblz}-3-\textsc{ipfv}\\
	\glt `It was raining frequently/all the time.' 	
	\ex \label{ex:fiebig:rainSete}
    \gll rɔ́p-\textbf{àn}-sàtɔ́-j\\
	\textsc{pst}-rain-\textsc{vblz}-sete-\textsc{ipfv}\\
	\glt `It was raining a lot/frequently.' 
\end{xlist}
\z

\ea Kipsigis \label{ex:fiebig:kip.rain}
\begin{xlist}
    \ex
    \gll rɔ̀ːp-tá\\
	rain-\textsc{sec}\\
	\glt `rain' (noun)
	\ex
    \gll kò-ròːp-\textbf{ɑ́n}\\
	\textsc{pst}-rain-\textsc{vblz}\\
	\glt `It rained.'
	\ex \label{ex:fiebig:rain-an-red}
    \gll ko-roːp-\textbf{ɑn}$\sim$ɑ$\sim$roːp.\textbf{ɑn}-i\\
	\textsc{pst}-rain-\textsc{vblz$\sim$lk}$\sim$rain.\textsc{vblz}\\ 
	\glt `It rained a lot/frequently.'
\end{xlist}
\z

\noindent The Pokot data in \REF{ex:fiebig:PokotRain}, which are identical to the Kipsigis data in \REF{ex:fiebig:kip.rain} (besides Kipsigis lacking the pluractional morpheme \textit{-sete}), show how the noun \textit{rɔ́p} `rain' can be verbalized by the suffix \textit{-an} (the b. examples respectively). To derive a pluractional meaning, the nominal stem plus its verbalizing suffix is reduplicated, as in \REF{ex:fiebig:rainRED}/\REF{ex:fiebig:rain-an-red} or modified by the suffix \textit{-sete}, as in \REF{ex:fiebig:rainSete}.\\
In theories like DM, lexical categories like nouns and verbs are decomposed into a categorizing head and a root. Morphemes like \textit{-an}, creating verbs from nouns can be analysed as such categorizing heads \citep{Marantz97}. In such an account, the data above show that the location of Kalenjin pluractionality is above little \textit{v} (the verbalizing/categorizing head) but below all other morphology. The structure in \figref{fig:fiebig:locPLAC} illustrates this observation.\footnote{Alternatively, \textit{-an} could verbalize an already nominalized structure \textit{n}P at this point: [\textit{v} [\textit{n}P [$\sqrt{\textsc{rop}}$ \textit{n}]] Under both alternatives, pluractional morphology applies after the categorizing head(s).} In \figref{fig:fiebig:locPLAC} and following structures, pluractionality (\textsc{plc}) is taken to be adjoined to the \textit{v}P and not project its own phrase. This is based on insights from \citet{Amato2018}, who treats pluractionality as an adverbial element (but see \citealt{Thornton2020} for an analysis where pluractionality is located on a head). At this point, I see no way to diagnose Kipsigis' and Pokots \textsc{plc} elements as either heads or adjuncts and therefore will remain agnostic about this distinction.


\begin{figure}
\begin{tikzpicture}[baseline=0pt] %baseline richtet Label und Beispiel zusammen aus
\tikzset{every tree node/.style={anchor=north},every node/.append style={align=center}}
\Tree 
[.\textit{v}P 
    [.\textit{v}P 
        [.$\surd$\\$\Rightarrow$rɔp ] \textit{v}\phantom{$\textsurd$}\\$\Rightarrow$-an ]  
    [.\textsc{plc}\\$\Rightarrow\sim$\textsc{red}/-sete ] 
] 
\end{tikzpicture}	
\caption{Location of pluractionality\label{fig:fiebig:locPLAC}}
\end{figure}

\noindent The ordering of the two pluractional morphemes in Pokot appears to be variable, with both orderings showing up, resulting in varying surface orders. These varying orders are indicated in the structures in \figref{fig:fiebig:varyB-structure} and \figref{fig:fiebig:vary-structure} below with the translations of the corresponding examples indicating their nuanced meaning differences (\REF{ex:fiebig:varyB} and \REF{ex:fiebig:vary} respectively).



\ea verb $\prec$ \textsc{red} $\prec$ \textit{-sete} \label{ex:fiebig:varyB} \\
\gll kìː-rál$\sim$àː$\sim$ràl-sàtá\\
to-cough$\sim$\textsc{lk}$\sim$\textbf{cough-sete}\\
\glt `to cough again and again' (maybe around 10 to 20 times)
\z   
    
\begin{figure}
        \begin{tikzpicture}[baseline=0pt] %baseline richtet Label und Beispiel zusammen aus
        \Tree 
        [.\textit{v}P 
            [.\textit{v}P 
                [.\textit{v}P $\sqrt{\textsc{root}}$+\textit{v} ]  
                [.\textsc{plc} \textsc{-red} ] 
            ] 
            [.\textsc{plc} -sete ] 
        ]
        \end{tikzpicture}
\caption{\label{fig:fiebig:varyB-structure}The structure of \REF{ex:fiebig:varyB}}
\end{figure}



\ea verb $\prec$ \textit{-sete} $\prec$ \textsc{red} \label{ex:fiebig:vary} \\
\gll kìː-rál-sàtà$\sim$ràl.sàtá\\
to-cough-sete$\sim$\textbf{cough-sete}\\
\glt `to cough again and again' (at least 10 to 20 times)
\z

\begin{figure}
\begin{tikzpicture}[baseline=0pt] %baseline richtet Label und Beispiel zusammen aus
        \Tree 
        [.\textit{v}P 
            [.\textit{v}P 
                [.\textit{v}P $\sqrt{\textsc{root}}$+\textit{v} ]  
                [.\textsc{plc} -sete ] 
            ] 
            [.\textsc{plc} \textsc{-red} ] 
        ]
        \end{tikzpicture}
\caption{The structure of \REF{ex:fiebig:vary}}\label{fig:fiebig:vary-structure}
\end{figure}

\noindent In \REF{ex:fiebig:varyB}, pluractionality by reduplication applies first. It therefore triggers the reduplication of the verbal stem (the root plus verbalizer) and the other pluractional morpheme \textit{-sete} applies after that. 
In contrast to that, in \REF{ex:fiebig:vary}, \textit{-sete} modifies the verb before reduplication applies. The reduplication inducing morpheme therefore triggers the reduplication of the verb plus \textit{-sete}. As the translations indicate, the meanings for this specific verb vary only slightly with their varying order: the meanings only seem to differ in the amount of times the event takes place. \REF{ex:fiebig:anotherone} provides another example where all possible orderings of reduplication and \textit{-sete} occur. 

\newpage
\ea 'to fold' in all possible orderings (Pokot)\label{ex:fiebig:anotherone}
\begin{xlist}
    \ex \label{ex:fiebig:anotheroneA}
    \gll kə̀-ɲàk̀ʊ̀m-àː-ɲàkʊ́m\\
    \textsc{pst}-fold-a-\textsc{red}\\
    \glt `to fold two or three times quickly' 
    \ex \label{ex:fiebig:anotheroneB}
    \gll kɛ́-ɲákʊ́m-sàtà\\
    \textsc{inf}-fold-sete\\
    \glt `to fold several times quickly'
    \ex
    \gll kɛ̀ː-ɲàkùm-àː-ɲàkùm-sàtà\\
    \textsc{inf}-fold-a-\textsc{red}-sata\\
    \glt `to fold again and again (more than \REF{ex:fiebig:anotheroneA} and \REF{ex:fiebig:anotheroneB})'
    \ex
    \gll kɛ̀ː-ɲàkʊ̀m-sàtà-ɲàkʊ̀m-sàtá\\
    \textsc{inf}-fold-sete-a-\textsc{red}\\
    \glt `to fold again and again (more than \REF{ex:fiebig:anotheroneA} and \REF{ex:fiebig:anotheroneB})'
\end{xlist}
\z

\noindent In this case, the ordering of reduplication and \textit{-sete} don't seem to make any difference, at least judging from the provided translations.\footnote{Often, the subtle meaning differences can only be pinned down by asking the speakers to provide the forms in imperative forms (e.g. ɲàkʊ́m-á-ɲákʊ́m-sátɛ̀-n-ə́ `fold again and again!'). The event is then acted out (folding a piece of paper two/three/fours times, folding it slowly/quickly and so on) and to the speaker is then asked if the order was fulfilled. Sometimes, meaning differences like in \REF{ex:fiebig:varyB} vs. \REF{ex:fiebig:vary} can be found in this way, but in many other instances, no clear meaning differences were identified yet.}\\
Overall, both pluractional morphemes (\textit{-sete} and reduplication) can be said to apply after little \textit{v}P (and not before it as, for example, analysed by \citealt{Thornton2019}) but below other, non-pluractional morphology. 


\section{Verbal number: participant and event plurality}\label{sec:fiebig:partNum}

Verbal number can be observed in two forms: participant and event number \citep{Corbett2000}. Examples \REF{ex:fiebig:part} and \REF{ex:fiebig:event} illustrate these two cases respectively.

\ea event number \label{ex:fiebig:event}
\begin{xlist} 
\ex Peter pets a dog once. \label{ex:fiebig:eventA}
\ex Peter pets the dog again and again. \label{ex:fiebig:eventB}
\end{xlist}
\z 


\ea participant number \label{ex:fiebig:part}
\begin{xlist} 
\ex Peter pets a dog. \label{ex:fiebig:partA}
\ex Peter pets several dogs. \label{ex:fiebig:partB}
\end{xlist}
\z

\noindent Event number differentiates if an event takes place once or several times. This distinction is illustrated by the contrast of \REF{ex:fiebig:eventA} and \REF{ex:fiebig:eventB}. The dog is petted once in \REF{ex:fiebig:eventA}, but several times in \REF{ex:fiebig:eventB}. Participant number indicates the number of participants. Examples \REF{ex:fiebig:partA} and \REF{ex:fiebig:partB} are different from each other in the number of participants, in this case objects. Both distinctions (event and participant number) are described as verbal number, since in both cases an event takes place several times: in \REF{ex:fiebig:eventB} the event happens repeatedly with the same participants and in \REF{ex:fiebig:partB} the event of petting happens repeatedly, distributed across different participants (dogs). In contrast to \citet{Corbett2000}, who classifies both of these phenomena as verbal number, other (especially semantic) literature considers event and participant number to be related but separate phenomena. In this literature event number is often called \textit{pluractionality} \citep[e.g.][]{Henderson2012, Wood2007}.
Coming back to the definition given in \REF{ex:fiebig:definition}, of course, I have to admit that the English examples in \REF{ex:fiebig:event} and \REF{ex:fiebig:part} are not really cases of verbal number, as they don't match the definition.
The differences of the a. and b. examples are not encoded by morphologically modifying the verb, but instead by adverbial constructions. Let us therefore have a look at a more classical example of verbal number. Example \REF{ex:fiebig:Mupun} from Mupun illustrates verbal number marking encoded directly on the verb.

\ea Verbal number in Mupun (Afro-Asiatic, Chadic) \citep[59-62]{Frajzyngier1993} \label{ex:fiebig:Mupun}
\begin{xlist}
\ex \label{ex:fiebig:hit.sgObj.sg}
\gll wu \framebox{cit} wur\\
 3\textsc{m.sg} hit.\textsc{pst.\textbf{sg}} \textsc{3m.\textbf{sg}}\\
\glt `He hit him.' 
\ex  \label{ex:fiebig:*hit.sgObj.pl}
\gll *wu \framebox{cit} mo\\
 3\textsc{m.sg} hit.\textsc{pst.\textbf{sg}} \textsc{3\textbf{pl}}\\ 
\glt `He hit them.' (participant number) 
\ex  \label{ex:fiebig:hit.sgObj.pl}
\gll wu \framebox{nas} mo\\
 3\textsc{m.sg} hit.\textsc{pst.\textbf{pl}} \textsc{3\textbf{pl}}\\
\glt `He hit them.' (participant number) 
\ex  \label{ex:fiebig:hit.plObj.sg}
\gll wu \framebox{nas} wur\\
 3\textsc{m.sg} hit.\textsc{pst.\textbf{pl}} \textsc{3m.\textbf{sg}}\\
\glt `He hit him \textbf{many times}.' (event number) 
\end{xlist}
\z

\noindent {Mupun exhibits a suppletive verb pair for the word `to hit': \textit{cit} and \textit{nas} for the singular and plural form respectively. The singular form \textit{cit} is used with a singular object and a singular event reading, as in \REF{ex:fiebig:hit.sgObj.sg}. The form cannot be used with a plural object, as the ungrammaticality of \REF{ex:fiebig:*hit.sgObj.pl} illustrates. To express the hitting of a plural object, the plural form \textit{nas} has to be used, as in \REF{ex:fiebig:hit.sgObj.pl}. The use of the plural form in this case is similar to the meaning difference described in \REF{ex:fiebig:part}: participant number. In contrast to that, in \REF{ex:fiebig:hit.plObj.sg} the plural verb in combination with a singular object expresses multiplicity of the same event happening to the same participant, similar to the meaning difference in \REF{ex:fiebig:event}: this is event number/pluractionality. The Mupun example therefore shows how both grammatical concepts related to verbal number are expressed by the very same marking strategy: both, participant number and pluractionality are expressed by the same suppletive form of a verb.}

Although not all languages conflate event and participant number in this way (Kipsigis and Pokot will turn out to not do so), data like \REF{ex:fiebig:Mupun} are taken as evidence that both flavours of verbal number should be analysed by the same mechanism. Therefore, the only two existing generative analyses of verbal number morphology (\citealt{Amato2018} and \citealt{Thornton2019}) both rely on the assumption that both flavours of verbal number trigger the same agreement on one head (v in Amato's and a \#-head in Thornton's account). In these analyses the head is then spelled out by pluractional morphology, independent of whether it was triggered by event or participant number.

But Pokot and Kipsigis pluractionals behave rather different from languages like Mupun: plural participant number does not trigger or require pluractional marking, only plural events do. 
In \REF{ex:fiebig:KipPl.OBJ} the plural object \textit{kwenik} `wood\textsubscript{\textsc{pl}}' is used with the unreduplicated form \textit{kiːabete} `I was splitting.' 

\ea Kipsigis 
\begin{xlist}
\ex 
\gll kiː-ɑ-bet-e kwen-deːt\\
\textsc{pst-1sg}-split-\textsc{ipfv} wood-\textsc{\textbf{sg}}\\
\glt `I was splitting a piece of wood.'
\ex \label{ex:fiebig:KipPl.OBJ}
\gll kiː-ɑ-bet-e kwen-ik\\
\textsc{pst-1sg}-split-\textsc{ipfv} wood-\textsc{\textbf{pl}}\\
\glt `I was splitting pieces of wood.' 
\ex \label{ex:fiebig:KipPLSG}
\gll kiː-ɑ-bet$\sim$ɑː$\sim$bet-i kwen-deːt\\
\textsc{pst-1sg}-split$\sim$\textsc{lk}$\sim$split\textsc{-ipfv} wood-\textsc{\textbf{sg}}\\
\glt `I was splitting one piece of wood \textbf{again and again}.' 
\ex \label{ex:fiebig:KipPLPL}
\gll kiː-ɑ-bet$\sim$ɑː$\sim$bet-i kwen-ik\\
\textsc{pst-1sg}-split$\sim$\textsc{lk}$\sim$split-\textsc{ipfv} wood-\textsc{\textbf{pl}}\\
\glt `I was splitting many pieces of wood \textbf{again and again}.’ 
\end{xlist}
\z

\noindent When the reduplicated (plural) form \textit{kiːɑbetɑːbeti} is used, a meaning of `the event takes place several times' is conveyed, with the number of the object being independent from that. This leads to the translation of several pieces of wood (plural object) being split several times (plural verb) in \REF{ex:fiebig:KipPLPL} and to the meaning of a single piece of wood (singular object) being split several times (plural verb) in \REF{ex:fiebig:KipPLSG}. Plural objects are possible with reduplicated and unreduplicated verbs in Kipsigis irrespective of aspect. As the following, aspectually unmarked examples illustrate, reduplicated and unreduplicated forms are grammatical without perfective marking. Furthermore, the reduplicated form still exhibits a multiple event reading, just as with the imperfective form.

\ea Kipsigis
\begin{xlist}
\ex 
\gll kɪ̀ː-bɛ́t knwɛ̀ːn-ɪ́k\\
\textsc{pst}-split firewood-\textsc{pl}\\
\glt `He split wood.'
\ex 
\gll kɪ̀ː-bɛ́t-àː-bɛ́t knwɛ̀ːn-ɪ́k\\
\textsc{pst}-split-\textsc{lk-red} wood-\textsc{pl}\\
\glt `He split wood again and again.'
\end{xlist}
\z

\noindent The same is true for both forms of Pokot pluractionality (only \textit{-sete} shown here). As \REF{ex:fiebig:sgpl} shows, the singular verb form \textit{kìtɛ́k} `he built' can be used freely with the plural object \textit{kɔ̀ːrín} `houses'. The verb doesn't have to be modified by a pluractional marker like reduplication or \textit{-sete}.

\ea Pokot
\begin{xlist}
\ex 
\gll kì-\fbox{tɛ́k} \fbox{kí} kɔ̀ɲíl ákɔ̀ŋə́\\
\textsc{pst}-\textbf{build} \textbf{house.\textsc{sg}} times one\\
\glt `There was one instance of him building a house.'
\ex \label{ex:fiebig:sgpl}
\gll kì-\fbox{tɛ́k} \fbox{kɔ̀ː-rín} kɔ̀ɲíl ákɔ̀ŋə́\\
\textsc{pst}-\textbf{build} house-\textbf{\textsc{pl}} times one\\ 
\glt `There was one instance of him building houses.' 
\ex \label{ex:fiebig:plpl}
\gll kì-\fbox{tɛ́k-sàtà} ɲíndə̀ \fbox{kɔ̀ː-rín}\\
\textsc{pst}-\textbf{build-sete} \textsc{pron.3sg} \textbf{house-\textsc{pl}}\\
\glt `He built houses several times.' 
\ex 
\gll kì-\fbox{tɛ́k-sàtà} ɲíndə̀ \fbox{kí}\\
\textsc{pst}-\textbf{build-sete} \textsc{pron.3sg} \textbf{house.\textsc{sg}}\\
\glt `He built a house very fast.'
\end{xlist}
\z

\noindent If, as in \REF{ex:fiebig:plpl}, a pluractional verb form is used, the sentence gets the meaning of `building several houses several times'. This, again, is in contrast to the Mupun data in \REF{ex:fiebig:Mupun}. In Mupun, the combination of plural verb and plural object doesn't receive a multiple event reading but instead only receives a multiple participant reading. The data above show that in both Kipsigis and Pokot reduplication and \textit{-sete} only express one of the two categories associated with verbal number: event number, but not participant number.

Kipsigis and Pokot employ a different marking strategy for participant number. Some Kalenjin verbs are suppletive for participant number. Both of these suppletive forms can co-occur freely with reduplication and \textit{-sete}.


\ea \label{ex:fiebig:suppl} Kipsigis \cite[7]{KouneliProp}
\begin{xlist}
\ex \label{ex:fiebig:suppl-sg}
\gll lɑbɑt-i làːkwɛ̀ːt. \\
run.\textsc{\textbf{sg}-ipfv} child.\textsc{nom} \\
\glt `The child is running.' 
\ex \label{ex:fiebig:suppl-pl}
\gll ruaj lɑ̀ːgôːk. \\
run.\textsc{\textbf{pl}.ipfv} children.\textsc{nom} \\
\glt `The children are running.' 
\ex \label{ex:fiebig:suppl-red}
\gll ruɑj$\sim$ruɑj lɑ̀ːgôːk.\\
run.\textbf{\textsc{pl}$\sim$run.\textsc{pl}} children.\textsc{nom}\\ 
\glt `The children are running \textbf{over and over again}.'
\end{xlist}
\z

\noindent The verb for `run' in Kipsigis is either spelled out as \textit{lɑbɑt} or \textit{ruaj}, depending on the number of the participant. In \REF{ex:fiebig:suppl-sg}, the form \textit{lɑbɑt} is chosen in the context of the singular participant/subject \textit{làːkwɛ̀ːt} `child'. In contrast to that, the form \textit{ruaj} is chosen in the context of plural subjects like \textit{lɑ̀ːgôːk} `children' in \REF{ex:fiebig:suppl-pl}. The reduplicative form \textit{ruɑj$\sim$ruɑ̀j} in \REF{ex:fiebig:suppl-red} shows that participant number (marked by suppletion) and event number (marked by reduplication) can freely cooccur in Kipsigis and are not in complementary distribution, as would be expected if they were spelling out the same head. 
\REF{ex:fiebig:supplPok} shows the very same for the two suppletive Pokot verbs \textit{rə̀p} `run' (singular) and \textit{rə́jɪ̀ɣ} `run (plural)'. Both suppletive forms (\textit{rə̀p} `run (singular)' and \textit{rə́jɪ̀ɣ} `run (plural)') can freely occur with both exponents of pluractionality: reduplication and \textit{-sete}.


\ea Pokot \label{ex:fiebig:supplPok}
\begin{xlist}
\ex \label{ex:fiebig:run-red}
\gll kì-rə̀p-à$\sim$rə́p\\
\textsc{pst}-run.\textsc{\textbf{sg}$\sim$lk}$\sim$run\\
\glt `S/he ran \textbf{repeatedly}.' 
\ex 
\gll kì-rə̀p-sɛ̀tə́\\
\textsc{pst}-run.\textsc{\textbf{sg}}-sete\\
\glt `S/he ran \textbf{repeatedly}.' (more often than \REF{ex:fiebig:run-red})
\ex \label{ex:fiebig:run.pl-red}
\gll kì-rə́jɪ́ɣ$\sim$à$\sim$rəjɪ̌ɣ\\
\textsc{pst}-run.\textsc{\textbf{pl}$\sim$lk}$\sim$run\\
\glt `They ran \textbf{repeatedly}.' 
\ex 
\gll kì-rə́jɪ́ɣ-sɛ̀tɔ̀\\
\textsc{pst}-run.\textsc{\textbf{pl}}-sete\\
\glt `They ran \textbf{repeatedly}.' (more often than \REF{ex:fiebig:run.pl-red})
\end{xlist}
\z


\noindent Kipsigis has an additional suffix to mark participant plurality in some cases. Just as with suppletion, this morpheme isn't in complementary distribution with event number marked by reduplication either, as \REF{ex:fiebig:toos} shows. The participant number suffix \textit{-toːs} can freely modify an already reduplicated verb like \textit{twɑl-twɑl} `jump over and over again' in \REF{ex:fiebig:toos-red}. 

\ea Kipsigis \cite[3]{KouneliProp}\label{ex:fiebig:toos}
\begin{xlist} 
\ex
\gll twɑl-toːs lɑ̀ːgôːk.\\
jump-\textsc{\textbf{pl}} children.\textsc{nom}\\
\glt `The children are jumping.' 
\ex \label{ex:fiebig:toos-red}
\gll twɑl$\sim$twɑl-to:s lɑ̀ːgô:k.\\
jump$\sim$\textbf{jump}-\textsc{\textbf{pl}} \textbf{children}.\textsc{nom}\\
\glt `The children are jumping over and over again.' 
\end{xlist}
\z

\noindent The above data indicate that in Kipsigis and Pokot, participant and event number are not two flavours of the same category as in Mupun. Instead, they are marked by different morphological strategies and are not in complementary distribution but can freely co-occur instead. Mupun and Kipsigis/Pokot therefore seem to form two extreme points of an continuum: on one end Mupun, where both categories are conflated fully and Kipsigis/Pokot on the other end, where both categories are fully distinct. It is of course expected that there are languages where the categories overlap in some aspects but not in others.

Although there is unfortunately not enough space to go into detail about the analyses of \citet{Amato2018} and \citet{Thornton2019}, my data show that Kipsigis and Pokot pose a problem for any analysis, that relies on the fact that event and participant number are located and spelled out on the same head (like the aforementioned ones do).


\section{Conclusion}

This paper provided the very first structural description of pluractionality in two Kalenjin languages: Kipsigis and Pokot. The semantic functions of pluractionality were investigated in detail for Pokot: the pluractional morphemes in Pokot express several semantic functions. One of the functions expressed by pluractionals cross-linguistically, habituality, is not expressed by pluractionality in Pokot, but rather by imperfectivity. \\
This paper further showed that in both languages, the pluractionality morphemes are added to the structure after the verbalizing head little \textit{v} but before any other, for example aspect related, morphology. In Pokot, where two pluractional morphemes can be observed, they co-occur in varying orders with respect to each other. Further, this paper investigated the relationship of event and participant number in both Kipsigis and Pokot, two categories that are often treated as flavours of the very same grammatical phenomenon: verbal number. It was possible to show that event and participant number are marked by different morphological strategies and are not in complementary distribution but can co-occur in both Kipsigis and Pokot. The results therefore suggest that participant and event number are two distinct categories at least in Kalenjin and cannot be accounted for by theories that rely on a conflation of both categories.\\
Further research needs to investigate if the distinction of both categories is a property of these two languages, the Kalenjin dialect cluster, or maybe their language family. Other questions for further investigations include the specific meaning differences of the two pluractional morphemes in Pokot and why Kipsigis has only one of them. 




\section*{Abbreviations}

Glosses follow the Leipzig glossing rules with the following additions:\\ the two Kalenjin pluractionality morphemes are glossed as \textsc{red} and \textit{-sete} respectively; \textsc{lk} = linker, \textsc{plc} = pluractionality, \textsc{sec} = secundary suffix, \textsc{vblz} = verbalizing head. Reduplicated morphemes are separated from their base with a tilde. Within reduplicants, morpheme boundaries that were present in the base reappear marked by a dot instead of the usual hyphen.


\section*{Acknowledgements}

I am grateful to Daisy Chepkopus Lochuliit, Enock Kirui, Henry Pkemoi, Hillary Mosonik and Scholah Loparatum for their valuable work as linguistic consultants. I am further grateful to the valuable feedback of the audience of ACAL 54 as well as the Leipzig university Morphology and Syntax colloquium, especially my supervisor Maria Kouneli. \\ I acknowledge the DFG grant FOR 5175: Zyklische Optimierung for funding.

\sloppy
\printbibliography[heading=subbibliography,notkeyword=this]
\end{document}
